% ICPE 2027 Research Track: 10 pages + references, ACM sigconf, double-anonymous.
% Build: tectonic main.tex   (or latexmk -pdf main.tex)
% Skeleton written 2026-11-14. Numbers come from phase2/*/RESULTS.md; every
% number must be re-checked against those files before submission.
% Draft only: font expansion off until the ACM fonts (texlive-fonts-extra) are installed.
\PassOptionsToPackage{expansion=false}{microtype}
\documentclass[sigconf,review,anonymous]{acmart}
\usepackage{booktabs}
\usepackage{graphicx}
\graphicspath{{../phase2/figures/}}
\usepackage{xcolor}
\newcommand{\todo}[1]{\textcolor{red}{[TODO: #1]}}
\newcommand{\uu}[1]{\,\textsf{#1}}

\setcopyright{none}
\acmConference[ICPE '27]{ACM/SPEC International Conference on Performance Engineering}{2027}{}
\settopmatter{printacmref=false}

\begin{document}

% TODO title: decide with supervisor
\title{The Hidden Cost of the Machine: How System Limits Slow Down AI Coding Agents}

\author{Anonymous Author(s)}
\affiliation{\institution{Anonymous Institution}\country{}}

\begin{abstract}
Cloud platforms run AI coding agents in containers with fixed CPU and
memory limits, from two vCPUs down to a quarter of a core. An agent waits
for a remote language model most of the time. It is therefore not known
where such limits start to cost, how much of the cost reaches the user,
or what the limits do to the agent. We ran the Claude Code agent 2,199
times on four cloud VMs with 2, 4 and 8 vCPUs, and recorded kernel traces
for almost every run. The runs cover ten kinds of system limits and
eleven real SWE-bench tasks at the container sizes that platforms offer,
in sessions of up to 24 minutes. We report three results. First, we map
the limits. Each limit has a breaking point and one of three shapes: a
smooth rise, a sudden cliff, or a crash. Some breaking points are the
same on every machine, and others grow with its size. Second, we measure
how much of a cost reaches the user. The model wait dilutes the cost, as
Amdahl's law predicts. At a quarter of a core, the agent's work on the
machine takes 3.9 to 4.8 times longer, and the user waits 1.8 to 2.6
times longer. Kernel traces show that a CPU quota adds waiting, not work,
so a simple formula predicts its cost. For disk limits, a replay of runs
with no limit predicts the time on the machine within a median 4.7\%,
where the simple formula is off by more than half. Third, we find
failures that only appear with live agents. The agent crashes when the
process limit is below the number of threads that its runtime keeps
alive, about 8 plus 2 per vCPU. When a test needs most of a memory limit,
the kernel kills the test and then, in 9 of 20 runs, the agent itself,
and the user gets no message. Giving the tools their own memory cgroup
removed these kills. A smaller model also did extra work when a test
printed a slow run time. We turn these results into sizing rules for
platform builders, and we release the dataset and tools.
\end{abstract}

\begin{CCSXML}
<ccs2012>
<concept><concept_id>10011007.10010940.10011003.10011002</concept_id>
<concept_desc>Software and its engineering~Software performance</concept_desc>
<concept_significance>500</concept_significance></concept>
<concept><concept_id>10010147.10010178</concept_id>
<concept_desc>Computing methodologies~Artificial intelligence</concept_desc>
<concept_significance>300</concept_significance></concept>
</ccs2012>
\end{CCSXML}
\ccsdesc[500]{Software and its engineering~Software performance}
\ccsdesc[300]{Computing methodologies~Artificial intelligence}

\keywords{AI coding agents, performance analysis, kernel tracing, cgroups, containers, LLM}

\maketitle

% ---------------------------------------------------------------------------
\section{Introduction}

AI coding agents such as Claude Code~\cite{claudecode} edit code, run
tests and install packages on their own. They usually run inside
containers, and the platform that hosts them decides how large each
container is. A GitHub-hosted CI runner for a
private repository gives 2 vCPUs and 8\,GB of memory~\cite{ghrunners}. A
default E2B agent sandbox gives 2 vCPUs and 512\,MiB~\cite{e2bsandbox}, a
Cloud Run service 1 vCPU and 512\,MiB by
default~\cite{cloudrun,cloudruncpu}, and the smallest AWS Fargate task a
quarter of a vCPU and 512\,MiB~\cite{fargate}. The Linux kernel
enforces these sizes with control groups (cgroups)~\cite{cgroupv2}. An
engineer who chooses a size, or who places many agents on one machine,
needs to know what each choice costs.

An agent spends much of its time waiting for the language model. In our
measurements this is 65 to 78\% of a run on small coding tasks and 42 to
66\% on real tasks, and a recent study of the same agent reports 71 to
98\%~\cite{yuan2026}. A small container therefore costs less than its
size suggests. How much less is not known, and neither is the size at
which the cost starts. Agent benchmarks report whether a task is solved
and what it cost in tokens, dollars or
turns~\cite{jimenez2024,yang2024}. They do not report what the system
did, so they cannot answer these questions. A recent study of agents in
sandboxes~\cite{agentcgroup} measured how much CPU and memory agents use
when nothing limits them. It applied memory limits only to replayed
memory traces, not to running agents, so it cannot say what a limit does
to an agent's run. A second study~\cite{choi2026} gave agent containers
2 to 8 CPU cores and slowed the model down on purpose. Extra cores helped
less when the model was slower. That study replayed recorded model
answers, however, so the agent could not change what it did. It also did
not test CPU shares below one core or other kinds of limits. These
studies show that work on the machine matters. What is missing is a
measurement of live agents under the limits that platforms apply: where
each limit starts to cost, how much of that cost reaches the user, and
what the agent does in response.

The model wait dilutes the machine's cost, but it does not remove it.
Two effects make the cost easy to misjudge. First, the time on the
machine is only a part of each run, so a large change there is a smaller
change in the total. This is Amdahl's law~\cite{amdahl1967} applied to a
slowdown. Figure~\ref{fig:timeline} shows it on one run: at a quarter of
a CPU core, the agent's work on the machine takes 3.9 times longer, but
the user waits only 1.8 times longer. Second, the model's answer time
varies by several seconds from run to run. With five runs per setting,
this variation hides any change in total run time below about 20\%. A
large cost is therefore visible but understated, and a moderate cost can
be hidden. For example, with a busy program next to the agent on the
same machine, the agent's work on the machine grows by 86\%. Amdahl's law
then predicts a 23\% longer run. We measured only 2 to 10\%, because the
model happened to answer about 3\,s faster in those runs.

\begin{figure}[t]
\centering
\includegraphics[width=\columnwidth]{ICPE-2027/figures/fig_timeline.pdf}
\Description{Two horizontal timelines of one run of the same coding task. With no limit the run takes 21 seconds: 15 seconds waiting for the model and 6 seconds of local work. At a quarter of a core it takes 39 seconds: the model wait is still 15 seconds, and the local work is 24 seconds.}
\caption{One typical run of the same coding task with no limit and at a
quarter of a CPU core (the runs with the median wall time, on a 4-vCPU
machine). The work on the machine takes 3.9 times longer. The model wait
does not change, so the run takes only 1.8 times longer.}
\label{fig:timeline}
\end{figure}

We measured this cost directly. We ran Claude Code on four cloud virtual
machines (VMs). Two VMs are of one type with 4 vCPUs. The other two have
2 and 8 vCPUs, and the 8-vCPU VM uses an AMD CPU instead of Intel. The
experiments cover ten kinds of limits: a CPU quota, soft and hard memory
limits, disk speed, network speed, a process limit, extra threads, a
busy neighbour, two limits at once, and up to 64 agents on one machine.
We also ran eleven real bug-fix tasks from SWE-bench
Verified~\cite{jimenez2024,swebenchverified} at the platform sizes above.
In total we collected 2,199 runs. A group of agents that share one
machine counts as one run. For 2,159 runs we recorded a kernel trace with
LTTng~\cite{desnoyers2006}. For the 2,067 runs on the three newer VMs, we
also recorded the agent's event log and all cgroup counters. The event
log is the agent's own record of each model call and each tool call, with
times. From it we split each run into four parts: start-up, waiting for
the model, running tools, and the agent program's own work. We call the
total run time \emph{wall time}. We call wall time minus the model wait
\emph{local time}. Local time is the part of a run that the machine
controls. Full tracing costs about 7\% more CPU time and adds no visible
wall time.

We answer five research questions:
\begin{itemize}
\item \textbf{RQ1} (anatomy): where does an agent's time go when nothing limits it?
\item \textbf{RQ2} (mechanisms): how does each limit change that picture?
\item \textbf{RQ3} (realistic use): do real platform sizes slow down real tasks?
\item \textbf{RQ4} (density): how many agents can share one machine?
\item \textbf{RQ5} (behaviour): does the agent change what it does because of the system?
\end{itemize}

Our main findings are:
\begin{itemize}
\item \textbf{Each limit has a breaking point and one of three shapes.}
  The breaking point is the first setting at which local time clearly
  grows. The shapes are a smooth rise (CPU quota, disk speed, threads), a
  cliff (memory), and a crash (process limit). On 2, 4 and 8 vCPUs, the
  breaking point of the CPU quota is 0.75 core, and the memory cliffs
  sit at the same sizes. The process limit that crashes the agent grows
  with the machine: 12, 16 and 24 on 2, 4 and 8 vCPUs. It is the number
  of threads that the agent's runtime keeps alive. Some of its thread
  pools have one thread per CPU, which gives about 8 plus 2 per vCPU on
  these three machines.
\item \textbf{The breaking points carry over to real tasks.} At 1 CPU or
  more, the short real tasks show no slowdown that we can detect. At 0.25
  CPU, every short task is still solved on every machine, but takes 1.35
  to 2.4 times longer. A long task whose memory need is close to 512\,MB
  takes 2.3 times longer at 1 CPU and 512\,MB, because of the memory
  limit and not the CPU.
\item \textbf{The model wait dilutes the cost, and runs with no limit
  predict it.} Let $f$ be the share of local time in a run with no limit,
  and let a limit make local time $s$ times longer. By Amdahl's law, the
  wall time grows by a factor of $1 + f\,(s-1)$. In our runs, $f$ is 0.20
  to 0.30 on small tasks and 0.17 to 0.73 on real tasks. Between a fifth
  and three quarters of a cost on the machine therefore reaches the
  user. The kernel traces show where $s$ comes from. Under a CPU quota of
  $q$ cores, the agent's threads use the same CPU time as before and
  spend the extra time waiting for the quota. For quotas below one core,
  $s = 1/q$ is therefore a good estimate. For disk limits it is not.
  There we replay the CPU and disk use of the runs with no limit, in
  steps of 0.2\,s, under the new limit. Nothing is fitted. On 89 held-out
  settings, the replay predicts the wall time within a median 6.9\% and
  local time within 7.1\%. On disk limits it predicts local time within
  4.7\%, where the simple estimate is off by 59\%.
\item \textbf{Limits cause failures that only live agents show.} When a
  test needs most of a memory limit, the kernel first kills the test,
  which is the largest process. In 9 of 20 runs it then kills the agent
  itself, and the user gets no message. Moving the tools into a child
  cgroup with their own memory limit removes these kills (0 of 20 runs)
  at no cost in time. On one coding task whose test prints its run time,
  a smaller model (Haiku) keeps working after the tests pass when the
  printed time is slow, and the run takes 2.1 to 2.6 times longer. One
  sentence in the prompt prevents this.
\item \textbf{CPU runs out before memory when agents share a machine.}
  For one coding task, throughput stops growing at about 4 agents per
  vCPU on all three machine sizes, whether the agents start together or
  arrive at random. At that point each agent needs about twice the CPU
  time that it needs alone. A capacity estimate based on one agent is
  therefore 2.3 to 2.5 times too high.
\end{itemize}

This paper makes four contributions:
\begin{enumerate}
\item A public dataset of 2,199 agent runs on four VMs, almost all with
  kernel traces, and the tool that collected it.
\item A map of ten limits and their breaking points on three machine
  sizes. It separates the effects that hold on every machine from the
  numbers that depend on the machine or the task.
\item A measurement of how much of a cost on the machine reaches the
  user, explained with Amdahl's law and kernel thread states, and a
  method that predicts the time under a limit from runs with no limit.
  We test it on held-out settings, machines and real tasks. We show
  where a simple estimate is enough (CPU quotas) and where the replay is
  needed (disk limits).
\item Failure modes that only appear with live agents, explained from the
  kernel trace. The process-limit crash is traced to the runtime's
  thread pools. The kill of the agent itself is traced to the kernel's
  choice of victim, and two simple, known fixes prevent it.
\end{enumerate}
Section~\ref{sec:design} describes the study design.
Sections~\ref{sec:rq1} to~\ref{sec:rq5} answer the research questions,
Section~\ref{sec:cross} compares the results across machines, and
Section~\ref{sec:implications} gives guidance for platform and agent
builders.

% ---------------------------------------------------------------------------
\section{Background and Related Work}
Our work connects five lines of research: the resource use of AI agents,
how agents react to their environment, resource limits and isolation,
performance variation and how to measure it, and kernel tracing. For each
line, we say how our study differs.

\emph{Resource use of AI agents.} Agent benchmarks score whether a task is
solved~\cite{jimenez2024,swebenchverified}. Agent platforms run each task
in a Docker sandbox with a shell and other tools~\cite{yang2024,wang2025}.
They do not report the CPU, the memory or the time that the sandbox gets.
Efficiency studies count turns (the steps of the agent's loop), tokens
and API cost~\cite{gao2026,lindenbauer2025,fan2025}.
SWE-Effi~\cite{fan2025} also splits run time into model time and local
work. It shows that the local share ranges from a small part of a run to
most of it, depending on the agent design. None of these studies limits
the machine.

Three recent studies measure the system side of agents.
AgentCgroup~\cite{agentcgroup} ran Claude Code in containers with no
limits. It found that tool calls take 20 to 35\% of the active time, and
that memory use comes in short spikes during tool calls. From the memory
peak of single agents, it estimated that memory, not CPU, limits how many
agents fit on a machine. Yuan et al.~\cite{yuan2026} ran Claude Code
against open models served on local GPUs. The model took 71 to 98\% of
the run time. They studied the serving side, where too many agents on one
GPU server break the reuse of cached model state. Choi et
al.~\cite{choi2026} replayed recorded model answers while they changed
the number of CPU cores of the tool containers and the speed of the
model. Extra cores helped less when the model was slower, and for coding
tasks the CPU or the disk, not memory, was the bottleneck. Our study
differs from these three in three ways. We run live agents under each
kind of limit, so the agent can react to the limit. We measure how many
live agents fit on a machine, and find that CPU fills first. We also
predict the time under a limit from runs with no limit.

\emph{Agents and their environment.} An agent learns about the machine
only through the output of its tools. SWE-agent showed that the design of
this interface changes how often an agent succeeds~\cite{yang2024}. Tool
output also makes up most of what the model reads: about 84\% of the
tokens of a SWE-agent turn in a pilot study~\cite{lindenbauer2025}.
Aegis~\cite{song2025} sorts failures caused by the environment into
three groups: exploration, exploitation and resource exhaustion, where
the resources are turns and tokens. It shows that agents carry on with
incomplete data when a tool returns only part of its output. It also
shows that changes to the tools, among them explicit reports of state
changes, reduce the failures in which the agent does not know the state
of its environment. A one-line reminder of the remaining turns also
changes how agents behave~\cite{gao2026}. Our RQ5 adds a class of failure
that comes from operating-system limits: a kill that the agent cannot
see, and a kill of the agent itself.

\emph{Resource limits and isolation.} The cgroup CPU quota is built on
the bandwidth control of the Linux scheduler~\cite{turner2010,cgroupv2}.
A cgroup may use a quota of CPU time in each period, and all its threads
stop when the quota is used up. Under this scheduler, containers can
receive less CPU than they reserve~\cite{liu2021}. For microservices, a
CPU quota at or just above the actual use of one service made the
slowest end-to-end requests five times slower~\cite{shetty2025}. We see
no such sharp jump for an agent, because the model wait dilutes each
stall.

Containers run CPU-bound code at close to native
speed~\cite{felter2015}. Lightweight sandboxes such as
Firecracker~\cite{agache2020} and gVisor run CPU-bound code at the same
speed as containers, but differ in their memory, disk and network
paths~\cite{anjali2020}. Firecracker still relies on the scheduler, the
memory manager and the cgroups of the host to share CPU and memory
between its users~\cite{agache2020}. This is the layer that we study.
Cgroups do not separate everything. Workloads in different cgroups still
share kernel locks, for example the lock of the kernel's memory
allocator~\cite{anjali2025}. The way a measurement is isolated also
matters. Duet benchmarking runs two versions of a program side by side
on one cloud machine, so that outside noise hits both~\cite{bulej2020}.
On the same cloud as ours, Docker isolation added false alarms to such
benchmarks, even with no noise, while cgroups with CPU pinning reduced
them~\cite{japke2025}.

Production platforms such as Google's Autopilot set CPU and memory
limits from past use, because a limit that is too low causes throttling
or out-of-memory kills~\cite{rzadca2020}. Autopilot does not model how a
limit changes the time a job takes. Sizeless predicts the run time of a
serverless function at other memory sizes from monitoring data of one
size~\cite{eismann2021}. Our estimate does the same for agents under a
CPU or disk limit, and it takes the model wait into account.

\emph{Performance variation and its measurement.} OS noise is the set of
short interruptions that the kernel and its background services cause.
It slows large parallel programs, because synchronization amplifies each
delay~\cite{petrini2003,tsafrir2005,ferreira2008}. Hoefler et
al.~\cite{hoefler2010} replayed measured noise traces in a simulator and
found that the pattern of the noise matters more than its average. An
agent shows the opposite of amplification: the long model wait dilutes
local delays in the wall time. Our replay drives a simple model with a
recorded run, like this noise simulation and like trace-based capacity
planning~\cite{gmach2007}.

Cloud machines add their own variation. Two VMs of one type can differ
more than one VM differs over time~\cite{leitner2016}. Hidden state at
the provider can make repeated runs depend on each
other~\cite{uta2020}. Performance data is rarely
normal~\cite{maricq2018}, and small details of the setup can flip a
conclusion~\cite{mytkowicz2009}. Several studies give rules for this
setting~\cite{hoefler2015,papadopoulos2021,laaber2019}: repeat runs,
randomize their order, run the baseline and the test on the same machine,
use non-parametric statistics, and report intervals. Our design follows
these rules (Section~\ref{sec:design}).

\emph{Kernel tracing and wait analysis.} LTTng~\cite{desnoyers2006}
records kernel events in per-CPU buffers at low cost. Wait analysis uses
the event that wakes a blocked thread to find why the thread waited.
From kernel events alone, it builds the critical path of a task, which
is the chain of waits that decides how long the task
takes~\cite{giraldeau2016}. Later work combines kernel and user-space
traces to find the cause of high latency~\cite{ezzati2018}, and builds
waiting dependency graphs of slow threads~\cite{ezzati2020}. These
methods assign waiting for a CPU to the threads that hold the CPU. Under
a cgroup quota, a thread can also wait while its core is idle. We
therefore add the quota as a separate cause of waiting.

% ---------------------------------------------------------------------------
\section{Study Design}\label{sec:design}

This section describes the machines and the agent, how we apply each
limit, the tasks, how we run and record each experiment, and how we
analyse the runs.

\subsection{Machines and agent}
All runs use Google Cloud VMs in one zone, with no swap and Ubuntu 24.04.
We call the first VM VM0 and the three later ones M0, M1 and M2. VM0 and
M0 are of type n2-standard-4. This type has 4 virtual CPUs (vCPUs) and
15\,GB of memory. Its CPU is an Intel Xeon Cascade Lake at 2.8\,GHz, and
the 4 vCPUs are two cores with two hardware threads each. M1
(n2-standard-2) has 2 vCPUs of the same CPU and 7\,GB of memory. M2
(n2d-standard-8) has 8 vCPUs of another vendor (AMD EPYC 7B13) and
31\,GB. VM0 produced 132 runs. It ran Linux 7.0.0-1011 and 7.0.0-1014,
with a disk that reads 205\,MB/s. M0, M1 and M2 run Linux 7.0.0-1017. M1
and M2 are copies of M0's disk, so they have the same kernel, tracer and
container images. All VMs use LTTng 2.15.1.

M0 holds two sets of runs. Set~1 (963 runs) used a disk that reads
287\,MB/s. For set~2 (623 runs) we moved M0 to a smaller disk that reads
187\,MB/s. M1 (230 runs) and M2 (251 runs) have the same 187\,MB/s disk.
Each set has its own runs with no limit, and we compare runs only within
one set. The reason is that cloud VMs of one type can differ from each
other~\cite{leitner2016,uta2020}. Disk I/O varies the most of all
resources, by an amount that depends on the
provider~\cite{leitner2016}.

The agent is Claude Code~\cite{claudecode}. We ran version 2.1.296 on
VM0, with the agent's default model at that time. On M0, set~1 and 68
runs of set~2 used version 2.1.322. The other 555 runs of set~2 and all
runs on M1 and M2 used version 2.1.327. On these three VMs the agent
uses the Opus model (\texttt{claude-opus-5-5}) unless we state Haiku
(\texttt{claude-haiku-5-5}) or Fable (\texttt{claude-fable-5-1}). The
agent runs without a user interface, with a fixed prompt and a fresh
session. On M0, M1 and M2, it also writes its event stream: one JSON line
for each model message and each tool call.

\subsection{Limits}
The agent and every process it starts run in their own cgroup. A cgroup
is a group of processes that the kernel limits and counts together. We
create an empty cgroup for each run, so that its counters start at zero.
A run sets a limit by writing one of the cgroup's control
files~\cite{cgroupv2}:
\begin{itemize}
\item \texttt{cpu.max}: the CPU quota, for example 0.25 of a core. The
  cgroup may use this share of CPU time in each period (100\,ms by
  default). Once the share is used up, the cgroup stops until the next
  period~\cite{turner2010,cgroupv2}.
\item \texttt{memory.high}: a soft memory limit. Above it, the kernel takes
  memory back from the cgroup (reclaim), but it never kills a process.
\item \texttt{memory.max}: a hard memory limit. When the kernel cannot take
  back enough memory, it kills a process in the cgroup.
\item \texttt{io.max}: the disk read speed.
\item \texttt{pids.max}: the number of processes and threads.
\end{itemize}
We write memory sizes in MB and GB. As in the kernel's control files,
these are binary units (1\,MB $= 2^{20}$ bytes), so 512\,MB is the same as
the 512\,MiB of the platforms above.

Three kinds of limits do not use a control file. Network limits use
\texttt{netem}, the Linux network emulator. We apply it in both
directions on a private network around the agent. A busy neighbour is a
program that loads the CPU, the memory or the disk from a separate cgroup
on the same VM. Extra threads and extra agents share the machine's CPUs
with no limit set. For real tasks, the agent works inside a Docker
container. The container's cgroup sits inside the agent's cgroup, so the
same limits apply to it.

\subsection{Tasks}
\emph{Small tasks} make each limit act on a known piece of work. Two of
them are coding tasks that the agent must solve. Task \texttt{t2} asks it
to fix four bugs across three modules. Task \texttt{t3} asks it to
implement functions from their descriptions until a test suite passes.
The other small tasks ask the agent to run a fixed workload, so that the
work itself does not change between runs:
\begin{itemize}
\item \texttt{t4} runs $W$ CPU-bound workers (programs that only compute)
  in parallel.
\item \texttt{t5} starts $N$ sub-agents that each run one worker.
\item \texttt{t6} reads a 1\,GB file several times (file memory).
\item \texttt{t7} runs a tool that allocates 512\,MB (process memory).
\item \texttt{t8} reads a 1\,GB file once (disk).
\end{itemize}
RQ5 uses four more tasks. The test of \texttt{t3} prints its own run
time. Task \texttt{t9} is \texttt{t3} without this line. Task
\texttt{t10} is a second coding task, with and without a printed run
time. Task \texttt{t11} has the same work as \texttt{t3} in three
variants: the test prints a fixed fast time, the test prints a fixed
slow time, or the prompt adds that run time is not a goal. Task
\texttt{t12} asks the agent to implement two functions. Its test suite
builds a 700\,MB buffer, and its test script sends the output through
\texttt{tee}. If the kernel kills the test, the output simply stops, and
the script still exits with 0. We chose this form because it is common.
A shell pipeline returns the exit status of its last command, so any
test script that pipes its output to a log or a filter hides a kill in
the same way. We also ran \texttt{t12} with two fixes that keep the
out-of-memory killer away from the agent (Section~\ref{sec:rq5}). Before
each run we restore the task's files to the same state. After the run we
check that the agent did not edit the test file.

\emph{Real tasks} are eleven issues from SWE-bench
Verified~\cite{jimenez2024,swebenchverified}, in six Python projects.
SWE-bench rates how long each fix would take a human. We call six tasks
\emph{short}: they are rated under 15 minutes (Django, Sympy, xarray,
pytest, Flask and Requests). We call four tasks \emph{long}: they are
rated 1 hour or more (Django, Sympy, pytest and xarray, with different
issues from the short ones). A fifth long task, a second Sympy issue
rated at more than 4 hours, gives sessions of many minutes. We build each
project at its starting version, with the dependency versions that the
SWE-bench harness fixes. We keep a task only if its target tests fail at
the start and pass with the reference fix in our environment. A run
solves the task when the target tests pass and the other tests that pass
with the reference fix still pass.

In set~1, the short tasks run in Docker at five sizes: no limit; 2 CPUs
and 8\,GB; 2 CPUs and 512\,MB; 1 CPU and 512\,MB; and 0.25 CPU and
512\,MB. The last four are the platform sizes of the introduction. The
short tasks also run with no container. The long tasks run with no limit
and at the two smallest sizes. Set~2 repeats the short tasks with no
limit and at the smallest size, and runs the fifth long task. M1 and M2
run three short tasks with no limit and at the two smallest sizes.

\subsection{Protocol}
Table~\ref{tab:data} lists the experiments. Each experiment changes one
setting step by step, from loose to tight. Set~1 of M0 has five runs per
step, with a few exceptions: three per step when we place many agents on
one machine, ten for sub-agents and for the tracing cost, and 15 for the
two main RQ5 groups. Set~2 of M0 adds runs around the breaking points and
for the main claims. It has 15 runs per CPU quota step, 10 per
process-limit step and per real-task setting, and 5 groups per step with
many agents. M1 and M2 repeat a core set of experiments with 5 to 10 runs
per step. On these two VMs, the numbers of threads and agents are scaled
to the number of vCPUs.

Before every run we empty the page cache, which is the kernel's
in-memory copy of recently read files. On VM0 we did this only for
memory and disk runs. Runs within an experiment are in random order,
with a fixed seed on M0, M1 and M2. Changes in the model's speed over the
day therefore cannot line up with one
step~\cite{maricq2018,hoefler2015}. A run is stopped after 15 minutes.
For real tasks the limit is 20 minutes for short tasks, 30 minutes for
long tasks, and 45 minutes for the fifth long task.

\begin{table}[t]
\caption{The dataset: 2,199 runs, 2,159 of them with a kernel trace (the
tracing-cost experiments include 40 runs without one). For many agents,
the count is groups of agents.}
\label{tab:data}
\resizebox{\columnwidth}{!}{%
\begin{tabular}{@{}lll rrrr@{}}
\toprule
& & & \multicolumn{2}{c}{M0} & & \\
\cmidrule(lr){4-5}
Experiment & Task & Steps & set 1 & set 2 & M1 & M2 \\
\midrule
CPU quota & t2, t3 & 2 to 0.25 core & 140 & 75 & 50 & 60 \\
Threads, sub-agents & t4, t5 & $C/2$ to $4C$ & 40 & -- & 20 & 20 \\
Memory (soft, hard, process) & t6, t7 & 2\,GB to 512\,MB & 90 & 30 & 35 & 35 \\
Disk read limit & t8 & 200 to 20\,MB/s & 30 & -- & 20 & 20 \\
Coding task: memory, disk & t3 & 512 to 64\,MB; 50 to 5\,MB/s & 65 & -- & -- & -- \\
Network & t3 & delay, loss, rate & 40 & -- & -- & -- \\
Process count & t3 & 64 to 8 & 70 & 60 & 45 & 50 \\
Neighbour, two limits & t3, t6, t8 & neighbours; pairs & 45 & 25 & -- & -- \\
Agents per machine & t3, Requests & 1 to 64 agents; staggered & 18 & 63 & 15 & 21 \\
Printed run time & t3, t9--t11 & 3 models & 140 & 105 & -- & -- \\
Silent test kill, two fixes & t12 & 768\,MB, none & -- & 80 & -- & -- \\
Real tasks, short & 6 or 3 tasks & 5 sizes; no container & 180 & 120 & 45 & 45 \\
Real tasks, long & 4 or 1 tasks & 3 platform sizes & 60 & 15 & -- & -- \\
Tracing cost & t3, Django & none, core, full & 45 & 50 & -- & -- \\
\midrule
Total & & & 963 & 623 & 230 & 251 \\
\midrule
VM0 & t2--t8 & 7 experiments & \multicolumn{4}{r}{132} \\
\bottomrule
\end{tabular}}
\end{table}

\subsection{What we record}
Each run keeps four views of the same events.
\begin{enumerate}
\item A kernel trace with LTTng. On M0, M1 and M2 it holds every kernel
  event and every system call, and each event names its process, thread,
  parent and namespace. VM0 used a smaller set: system calls, scheduling,
  disk, network, interrupts, power and memory reclaim.
\item The agent's event stream, with each line time-stamped when it
  arrives (not on VM0).
\item Every cgroup counter (\texttt{cpu.stat}, \texttt{memory.*},
  \texttt{io.stat}, \texttt{pids.*}), read before and after the run and 5
  times per second during it. VM0 keeps only \texttt{cpu.stat} before and
  after the run, and the agent's final answer.
\item The versions of the kernel, tracer, agent and collection tool, and
  shared timestamps that line up the three other views.
\end{enumerate}

\subsection{Splitting a run}
From the event stream we split the wall time of each run into parts. A
\emph{model turn} starts when the agent hands work to the model, at the
start of the session or with a tool result. It ends with the last block
of the model's answer. A \emph{tool call} starts when the model asks for
a tool and ends with the tool's result. We assign every 10\,ms of the
run to one part:
\begin{itemize}
\item \emph{start-up}: before the first model turn;
\item \emph{model wait}: a model turn is open;
\item \emph{tool time}: only a tool call is open;
\item the agent's \emph{own work}: neither is open;
\item a short \emph{shut-down} at the end.
\end{itemize}
We take model turns from the event stream and not from network traffic.
A rule based on traffic gives wrong turns on slow networks, with small
models and with sub-agents.

\emph{Local time} is the wall time minus the model wait. It is the part of
a run that the machine controls. Some work on the machine runs while a
model answer is still open. This work counts as model time, so local time
is a lower bound. We used the trace to measure how much work is hidden in
this way. For every process the run started, apart from the agent itself,
we took the part of its lifetime that falls inside an open model answer.
We did this for 1,836 traced single-agent runs on M0, M1 and M2. These
are all experiments except the tracing-cost and long-session runs, and
without the runs in which the agent aborted before its first model turn.
The hidden work is 0.25\% of local time in the median run (90th
percentile 5\%). It is 0.2\% on small tasks and 1.6\% on real tasks. On
real tasks it is mostly a \texttt{git} call that the agent makes at the
start of a session, while the first model answer arrives. In the Django
repository this call takes up to 10\,s. The hidden share does not grow
under limits: under a CPU quota it is 0.08\% (median; 90th percentile
0.7\%). Moving the hidden work into local time changes no conclusion.
The error of the wall-time formula of Section~\ref{sec:rq2} stays at 3.5
to 3.6\% (median over 75 settings).

From the trace we also split each thread's time into three states:
running, waiting for a CPU, and sleeping. Waiting for a CPU means waiting
for the CPU quota or for a busy core. Sleeping means waiting for the
disk, for memory reclaim, or for another cause. This follows the method
of kernel wait analysis~\cite{giraldeau2016,ezzati2020}. That method
assigns waiting for a CPU to the threads that hold the CPU. Under a
quota, a thread can also wait while its core is idle, so we count this
case separately. We check this split against the kernel's own counter in
Section~\ref{sec:rq2}.

\subsection{Statistics}
We call one combination of task, model and limit value, on one machine
and set, a \emph{setting}. We compare each setting with the same task
and model with no limit. Where an experiment has no runs with no limit,
we compare with its loosest step. The \emph{slowdown} is the median time
in the setting divided by the median time with no limit. Medians are not
pulled by the few very slow model
answers~\cite{maricq2018,hoefler2015}. We give a 95\% bootstrap interval
for each slowdown: we draw the runs again with replacement 2,000 times
and take the middle 95\% of the resulting ratios.

Going from loose to tight, the \emph{breaking point} is the first step
whose whole interval for local slowdown lies above 1.10. In other words,
it is the first step where we are confident that local time grows by
more than 10\%. We use local time because the model's answer time varies
a lot. With five runs, a change in wall time below about 20\% cannot be
told apart from this variation, while a change in local time of about
10\% can. The 10\% threshold is a choice. It is close to the smallest
slowdown that most cloud benchmark setups can detect when the baseline
and the test run on the same instances in random, interleaved
order~\cite{laaber2019}. That study used 10 to 20 instances, and with
one instance detection was much weaker. We use one VM per set. We
compensate with more runs per step and by repeating the main experiments
on three VMs.

We checked the breaking-point rule in four ways, with the runs we have.
\begin{itemize}
\item \emph{Other thresholds.} We computed every breaking point at 5\%
  and 20\% as well. On local time, the breaking points of the CPU quota,
  the soft memory limit and the process limit do not move on any
  machine. Of all 42 breaking points, 24 stay the same, 15 move by one
  step and 3 move by more, mostly on wall time.
\item \emph{Means instead of medians.} We repeated every slowdown and
  every breaking point with means. The local slowdown from means differs
  from the one from medians by 1.1\% (median over all settings; 4.5\% at
  the 90th percentile). Of the 42 breaking points, 37 are the same with
  means. On local time, those of the CPU quota, the soft memory limit
  and the process limit are the same on every machine. Extra threads
  start to cost one step earlier on two machines. Where a mean and a
  median lead to different statements, we say so in the text.
\item \emph{A/A test}~\cite{laaber2019}. For every setting with no limit
  and 10 or 15 runs, we split the runs into two random halves, 400
  times. The rule called one half slower in 0.1\% of the splits, for
  local and for wall time.
\item \emph{Five runs.} We drew five runs per side from the settings
  with 10 to 15 runs, as in set~1, again 400 times per setting. The 95\%
  interval from five runs held the value from all runs in 99\% of the
  draws. The breaking-point verdict agreed with the one from all runs in
  90\%. Every disagreement was a missed breaking point close to the
  threshold, such as 0.75 core, where local time grows by 28\%. None was
  a false one. With five runs, a breaking point can therefore appear one
  step tighter than it is, but not looser.
\end{itemize}

To compare two groups of runs we use a one-sided rank-sum test. The test
is exact: it uses every way to split the runs into two groups. For the
larger RQ5 groups we use 200,000 random splits instead. Where we run many
tests, we apply a Holm correction, which keeps the chance of any false
finding among them at 5\%. As effect size we give Cliff's $\delta$. It is
the share of pairs of runs in which the first group is slower, minus the
share in which it is faster, and goes from $-1$ to 1. For rates of
solved tasks, kills and crashes we give a 95\% Wilson interval, and we
compare two rates with a one-sided Fisher exact test. To compare two
predictions on the same settings we use a two-sided sign test.

We call a result \emph{confirmatory} when the runs that test it were
designed after the claim was made from other runs. Set~2 of M0, M1 and M2
test the claims of set~1. The follow-up runs at process limits 14 to 8
(M1) and 64 to 28 (M2) test the process-limit rule of
Section~\ref{sec:cross}. The runs with the two out-of-memory fixes test
the cause that we proposed for the agent being killed in \texttt{t12}.
The limit of four agents per vCPU was seen first in set~1. It was then
tested on set~2, M1 and M2, and with agents that arrive at random.
Results seen first in the new runs are exploratory. Examples are the
agent being killed in \texttt{t12} and the long-session results.

\subsection{Checks on the measurement}
\emph{Tracing cost.} We ran the coding task \texttt{t3} ten times each
with no trace, with a core profile (all system calls and the main
scheduling, disk, network and memory events), and with the full profile.
We also ran a real Django task five times with each of the three. Wall
time did not change. On the coding task the full trace added about 7\%
CPU time (0.4\,s). On the real task we found no cost.

A trace could cost more when CPU is short. In set~2 we therefore
repeated the test, with no trace against the full trace, in three
settings: at 0.25 core (10 runs each), at a 192\,MB memory limit (10 runs
each) and with 8 agents at once (5 groups each). Wall time did not change
in any of them. The ratio of wall times is 1.00 [0.94--1.07] at 0.25
core, 1.01 [0.88--1.11] at 192\,MB and 1.08 [0.93--1.25] with 8 agents.
The CPU time of the agent's cgroup grows by 2 to 8\%. No trace lost an
event. Every ratio we report compares two traced runs, so most of this
cost cancels.

The tracer's daemon, which writes the trace buffers to disk, runs as a
system service outside the agent's cgroup. Its CPU time is therefore not
in these numbers. It writes to the same disk as the agent, so its writes
share the disk with the agent, but they do not count against the agent's
disk limit. Other studies report a higher cost for full kernel tracing:
up to 42\% on a disk-heavy program when the trace is written to the same
disk~\cite{ezzati2020}. A likely reason for the difference is that our
agent spends most of a run waiting for the model, and the machine
produces few kernel events during that time.

\emph{Trace against counters.} The CPU time measured from the trace agrees
with the kernel's own counter (\texttt{cpu.stat}). The difference is
0.8\% in the median run and 2.0\% at the 90th percentile, over the 1,952
runs that we analysed from the trace.

% ---------------------------------------------------------------------------
\section{RQ1: Where Does the Time Go?}\label{sec:rq1}

\begin{figure}[t]
\centering
\includegraphics[width=\columnwidth]{ICPE-2027/figures/fig_split.pdf}
\Description{Horizontal stacked bars for two small coding tasks (four bars, three models) and six real tasks, showing the share of wall time in start-up, model wait, tools and other. The model wait is the largest part in all bars; tools are larger on real tasks.}
\caption{Where the time goes with no limit: mean share of wall time per
part (Claude Code; the six short real tasks with Opus, in Docker). Median
wall time on the right. The shares in the text are the same means.}
\label{fig:split}
\end{figure}

Before we apply any limit, we look at where an agent's time goes when
nothing limits it. Figure~\ref{fig:split} splits runs with no limit into
the parts of Section~\ref{sec:design}. It uses the 83 runs of set~1 with
no limit on the two coding tasks and the six short real tasks.

On the small coding tasks, the agent waits for the model for 72 to 78\% of
the time with Opus, 65\% with Haiku and 75\% with Fable (means). Tools
take 12 to 24\% of the time and start-up 7 to 9\%. The agent's own work
between events takes about 0.1\%, and shut-down about 2\%. With no limit,
the system adds almost no waiting: the agent's processes wait for a CPU
for only about 0.1 to 0.2\,s per run.

Real tasks move time from the model to the tools, because real test
suites do more work. On the six short real tasks, tools take 20 to 40\%
of the time (the most on xarray and Flask) and the model 42 to 66\%.
Start-up is a fixed cost: 2.6\,s without a container and 3.8\,s in Docker
(medians). The container therefore adds 1.2\,s to every start. The rest
of a run is not clearly slower in a container.

The largest source of spread between runs is the agent's own choices.
The slowest short run with no limit took 85\,s. The other runs of the
same task (Requests), with and without a container, took 20 to 27\,s. In
the slow run the agent chose to run the project's whole test suite. For
this reason, before we blame a limit for a change in tool time, we
compare the \emph{same} test command across settings
(Section~\ref{sec:rq3}).

\textbf{Answer to RQ1.} With no limit, the agent waits for the model for
65 to 78\% of a run on small tasks, and for 42 to 66\% on real tasks,
which run heavier tests. The machine works for the rest: a fifth to a
third of the run on small tasks, and up to more than half on real ones.
Any cost on the machine is therefore diluted by the model wait. The next
section measures by how much.

% ---------------------------------------------------------------------------
\section{RQ2: How Does Each Limit Change the Picture?}\label{sec:rq2}

\begin{figure*}[t]
\centering
\includegraphics[width=\textwidth]{ICPE-2027/figures/fig_limits.pdf}
\Description{Five line charts of wall and local slowdown against limit level for CPU quota, disk read limit, workers, soft memory limit and disk limit on the coding task, each with its breaking point marked. A sixth chart shows the share of runs in which the agent aborts against the process limit: it jumps from 0 to 100 percent at 12 on 2 vCPUs, at 16 on 4 vCPUs and at 24 on 8 vCPUs.}
\caption{The three shapes. First five panels: slowdown against the limit
level, from loose (left) to tight (right), on M0 with Opus. Wall is what
the user waits; local is wall minus model wait. Medians of 5 runs with
95\% bootstrap intervals; the dotted line marks the breaking point. Last
panel: share of runs in which the agent aborts under a process limit, on
the three machine sizes (5 to 15 runs per point). On 2 vCPUs, 3 of 5 runs
abort at a limit of 10, between 5 of 5 at 12 and at 8; we cannot explain
this dip with five runs.}
\label{fig:limits}
\end{figure*}

\begin{table}[t]
\caption{Each limit at a glance (Opus, local slowdown against no limit).
Breaking point: going from loose to tight, the first step at which local
time grows by more than 10\%. For the neighbour, the agent itself has no
limit, and we give the cost of each kind of load.}
\label{tab:limits}
\resizebox{\columnwidth}{!}{%
\begin{tabular}{@{}llll@{}}
\toprule
Limit & Breaking point & Tightest step & Part that grows \\
\midrule
CPU quota & 0.75 core (+35\%) & 0.25: 3.9$\times$ (wall 2.0$\times$) & tools, start-up \\
Threads on 4 vCPUs & 8 workers (+60\%) & 16: 2.4$\times$ & tools (CPU wait) \\
Disk read speed & 150\,MB/s (+43\%) & 20\,MB/s: 9.2$\times$ & tools \\
Disk, coding task & 50\,MB/s (+26\%) & 5\,MB/s: 5.3$\times$ & start-up \\
File memory, soft/hard & 1152 / 1024\,MB & 3.4$\times$ in one step & tools (re-reads) \\
Tool memory (512\,MB) & 640\,MB & tool killed, 5 of 5 & model (explains kill) \\
Whole-agent memory & 128\,MB & 15$\times$, then killed & start-up, reclaim \\
Process count & 17 (+55\%) & $\leq$16: most abort in 3\,s & tools \\
Network & none & 256\,kbit/s: wall +29\% & model only \\
Busy neighbour & CPU load: +86\% & memory load: +62\% & tools, start-up \\
\bottomrule
\end{tabular}}
\end{table}

We now apply one limit at a time and tighten it step by step.
Table~\ref{tab:limits} summarizes every limit, and Figure~\ref{fig:limits}
shows six of them. Tasks \texttt{t4} to \texttt{t8} are fixed workloads
that the agent only starts. For them, the response to a limit is known
kernel behaviour. We use them as calibration: they show what each limit
does to a known piece of work. Four results are specific to agents:
where the cost lands in an agent's run, the process-limit crash, the
kills, and how much of the cost reaches the user. The limits follow
three shapes.

\emph{Smooth.} A CPU quota, a disk limit and extra threads raise the cost
step by step, roughly in proportion to the cut.
\begin{itemize}
\item \emph{CPU quota.} At 0.75 core, the local time of the coding task
  \texttt{t3} grows by 35\%. At 0.5 core it grows by 84\%, and at 0.25
  core it is 3.9 times longer. At 0.25 core, the extra time is in the
  tools (13\,s more) and in start-up (3.6\,s more). The trace shows that
  the agent's processes wait 26\,s for the CPU quota.
\item \emph{Disk limit.} A disk limit acts in the same way on a file
  read: at 20\,MB/s, local time is 9.2 times longer. On the coding task,
  a disk limit slows start-up instead (21\,s more at 5\,MB/s). The agent
  program itself is read from disk when it starts, because we empty the
  page cache before each run. On a machine where the agent's files stay
  cached between runs, this cost would be smaller.
\item \emph{Extra threads.} More workers than the four vCPUs add waiting
  for a CPU: 7.6\,s with 8 workers and 46\,s with 16.
\end{itemize}

\emph{Cliff.} A memory limit has no effect while the data that a task
needs still fits. Once the data no longer fits, the cost jumps in one
step. The task \texttt{t6} reads a 1\,GB file several times. Under a soft
limit (\texttt{memory.high}), it runs at full speed down to 1344\,MB. At
1152\,MB it is 3.4 times slower, because every pass reads the file from
disk again. A hard limit (\texttt{memory.max}) gives the same cliff for
file memory, at 1024\,MB, and never kills. File data in memory is only a
copy of the disk, so the kernel can always drop it and read it again. A
program's own memory has no such copy, and our VMs have no swap, so the
kernel can only kill the program. The tool of \texttt{t7} needs 512\,MB.
It runs at 704\,MB and is killed at 640\,MB in all five runs. A limit on
the whole agent on the coding task costs 11\% at 192\,MB. At 128\,MB the
run takes 15 times longer, spent on reclaim and disk reads, before it is
killed. Below that, the agent dies during start-up.

\emph{Crash.} A process limit has no cost until the agent crashes. In
set~1, some of the agent's attempts to start a process fail from a limit
of 20 down, but the shell retries them. At 17, the run is 55\% slower. At
16, four of five runs abort within about three seconds, and at 12 all
five do. In set~2, with 10 runs per step, the edge is the same. All 10
runs pass at 18. At 17, all 10 solve the task, but one agent aborts after
solving it. At 16, all 10 abort. The agent thus almost always crashes
below the edge and rarely above it. Above the edge, 3 of 55 runs in
set~1 and 1 of 50 in set~2 aborted. The last panel of
Figure~\ref{fig:limits} shows this edge, and that it moves with the
number of vCPUs (Section~\ref{sec:cross}).

Three cases do not fit these shapes.

\emph{Network.} The network is the only limit that leaves local time
without a breaking point. Delay, packet loss and low bandwidth add up to
6\,s of model time and nothing clear on the machine.

\emph{Busy neighbour.} A busy neighbour raises the agent's local time by
86\% (CPU load) or 62\% (memory load), although the agent itself has no
limit. Amdahl's law predicts that wall time grows by 16 to 23\%. We
measured only 2 to 10\%, because the model answered 2 to 3\,s faster in
these runs than in the runs with no neighbour.

\emph{Two limits at once.} Two limits together can cost less or more
than their separate costs combined. We compare each pair with two
references: the sum of the two extra costs, and their product. Each pair
was measured against its own runs with no limit, so the single-limit
values below differ slightly from those above.
\begin{itemize}
\item \emph{CPU and disk cost less than their sum.} On \texttt{t3} in
  set~2, a CPU quota of 0.5 core alone makes local time 1.92 times
  longer, and a disk limit of 20\,MB/s alone 1.84 times. Together they
  make it 2.53 times longer, less than both the sum (2.76) and the
  product (3.53). The quota slows the CPU-heavy moments of the run and
  the disk limit slows the reads. While the run waits for one, it does
  not also wait for the other.
\item \emph{CPU and memory cost a little more.} In set~1, a CPU quota of
  0.5 core (1.82) and a memory limit of 256\,MB (1.04) cost 2.04
  together, a little more than both the sum (1.86) and the product
  (1.89).
\item \emph{A process limit adds nothing to a memory limit.} With a
  192\,MB memory limit, local time is 1.23 times longer alone and 1.22
  with a process limit of 20 added.
\item \emph{Memory and disk interact most.} A soft memory limit of
  768\,MB alone slows the task \texttt{t6} from 15.5\,s to 34\,s. With a
  50\,MB/s disk limit added, it takes 117\,s, and its processes wait
  96\,s for the disk. The memory limit turns reads from memory into
  reads from disk, and the disk limit then slows these reads down. We
  did not run the disk limit alone on this task.
\end{itemize}

\subsection{What the kernel trace shows}

\begin{figure}[t]
\centering
\includegraphics[width=\columnwidth]{ICPE-2027/figures/fig_states.pdf}
\Description{Four panels of stacked horizontal bars of thread time by state. Under a CPU quota, running time stays near 6 seconds while waiting for the quota grows to 26 seconds at a quarter of a core. With more workers than CPUs, waiting for a busy core grows to 46 seconds at 16 workers. Under a disk read limit and under a soft memory limit, only waiting for the disk grows.}
\caption{Where the threads of the agent's cgroup spend their time, from
the kernel trace (M0, set~1, medians). Thread time adds up over all
threads, so it can be longer than the wall time. A CPU quota, extra
threads and a disk limit each grow one kind of waiting and leave the
running time unchanged. A memory limit also adds running time, because
the file is read again.}
\label{fig:states}
\end{figure}

The slowdowns above say how much each limit costs. The kernel trace says
why. For every thread in the agent's cgroup, the trace gives the time it
ran, the time it was ready but waited for a CPU, and the time it slept
(Section~\ref{sec:design}). Figure~\ref{fig:states} shows these states for
four limits.

\emph{A CPU quota adds waiting, not work.} On the coding task, the
threads run for 6.3\,s with no limit and for 6.7\,s at 0.25 core. Their
time waiting for the quota grows from 0.1\,s with no limit to 2.5\,s at
0.75 core, 7.1\,s at 0.5 core and 26.4\,s at 0.25 core. They never wait
for a busy core, because the machine has free CPUs. The agent therefore
does the same work under the quota. This is why the work measured with
no limit can predict the time under a quota.

\emph{Too many threads look different.} With 16 workers on 4 vCPUs, the
threads wait 45.9\,s for a busy core and 0.1\,s for the quota. Standard
wait analysis reports both cases as ``waiting for a CPU''. We separate
them with one rule. A thread that is ready to run while a CPU is idle is
held back by the quota. A ready thread with no idle CPU waits for a busy
core. We checked this rule against the kernel's own counter of throttled
time (\texttt{cpu.stat}). In the 275 runs with more than 1\,s of
throttling, the quota wait from the trace is 1.10 times the counter
(median; 1.03 to 1.33 for the middle 80\% of runs). It is above 1 because
several threads can wait at the same moment. An operator who sees CPU
waiting can thus tell a quota that is too small from a machine that is
too full.

\emph{Disk and memory limits grow disk waiting.} Under a read limit of
20\,MB/s, the file read waits 57\,s for the disk, against 4.5\,s with no
limit. Under a soft memory limit of 1152\,MB, the same kind of waiting
grows from 4.5 to 19.5\,s. The CPU time also grows, from 3.2 to 5.4\,s,
because every pass reads the file again. The cliff of the memory limit is
thus a disk cost.

\emph{A busy neighbour raises the agent's CPU time.} With four CPU-bound
programs next to the agent, its threads wait only 3.0\,s for a busy core.
But they run for 10.4\,s instead of 6.2\,s to do the same work. The
neighbour makes each unit of work cost more CPU time. The likely cause is
that the vCPUs are hardware threads that share cores and caches. We come
back to this effect in Section~\ref{sec:rq4}.

\subsection{How much reaches the user: Amdahl's law and a prediction}

In the first five panels of Figure~\ref{fig:limits}, the wall slowdown is
much smaller than the local slowdown. Amdahl's law~\cite{amdahl1967}
explains this. Split a run with no limit into model time and local time,
and let $f$ be the share of local time. If a limit makes local time $s$
times longer and leaves the model's answer time unchanged, the wall
slowdown is
\begin{equation}
\text{wall slowdown} = 1 + f\,(s - 1).
\label{eq:hidden}
\end{equation}
The formula is not new. Our contribution is to measure $f$ and $s$ for
agents and to test where its assumption fails. The formula alone does
not predict a cost, because $s$ is known only after a run under the
limit.

\emph{How large is $f$?} In runs with no limit, $f$ is 0.20 to 0.30 on
the small coding tasks and 0.17 to 0.73 on the real tasks (medians per
task). This range is wider than the shares of Section~\ref{sec:rq1},
because it covers all sets, all machines and the long tasks. The highest
value belongs to the long Sympy task of Section~\ref{sec:rq3}. Between a fifth and three quarters of a cost on
the machine therefore reaches the user. A large cost stays visible: at
0.25 core, a local slowdown of 3.9 gives a measured wall slowdown of 2.0.
A moderate cost may not. On the small tasks, a local slowdown of 1.5
becomes 1.10 to 1.15. This is below what five runs can tell apart from
the variation of the model's answer time. OS noise in large parallel
programs behaves in the opposite way: there, synchronization amplifies
small local delays~\cite{hoefler2010}.

\emph{Does the formula hold?} We tested Equation~\ref{eq:hidden} with the
measured $s$ on the 68 limited settings of set~1. It is within 10\% of
the measured wall slowdown in 52 of them. In 45 of the 68 settings the
model time changed by less than 2\,s, and the formula is within 10\% in
44 of these. Of the 16 misses, 15 are settings where the model time
changed by 2\,s or more, and the miss has the direction of that change.
Examples are slower answers under a network limit and extra model turns
(Section~\ref{sec:rq5}). The error of the formula therefore measures the
part of a slowdown that comes from the model side.

\begin{figure}[t]
\centering
\includegraphics[width=\columnwidth]{ICPE-2027/figures/fig_hidden.pdf}
\Description{Scatter plot of measured against predicted wall slowdown on logarithmic axes, one point per setting, for four sets of runs. Most points lie inside a band of plus or minus 10 percent around the diagonal. Almost all points far above the diagonal are hollow, meaning the model time changed.}
\caption{Wall slowdown predicted by the replay against the measured one.
Each point is one setting with a CPU quota or a read limit. Hollow points
are settings where the model time changed by 2\,s or more, that is, where
the agent did something different.}
\label{fig:hidden}
\end{figure}

\begin{table}[t]
\caption{Predicting the time under a CPU quota or a read limit from the
runs with no limit. 89 held-out settings (small and real tasks, three
models, all four sets of runs). Absolute error of the wall time (median
with 95\% bootstrap interval, mean, settings within 10\%) and of local
time, the part that a limit changes (median, 90th percentile).}
\label{tab:predict}
\resizebox{\columnwidth}{!}{%
\begin{tabular}{@{}lrrrrr@{}}
\toprule
& \multicolumn{3}{c}{Wall time} & \multicolumn{2}{c}{Local time} \\
\cmidrule(lr){2-4}\cmidrule(l){5-6}
Prediction & Median & Mean & $\leq$10\% & Median & 90th pct. \\
\midrule
\textbf{Replay of 0.2\,s slices} & \textbf{6.9\%} [5.2--8.2] & \textbf{11.6\%} & \textbf{61} & \textbf{7.1\%} & \textbf{40\%} \\
Same rule on the whole local part & 8.0\% & 12.5\% & 56 & 11.8\% & 51\% \\
Eq.~\ref{eq:hidden} with $s = \max(1, 1/q)$ & 9.4\% [6.7--11.9] & 17.7\% & 47 & 8.7\% & 68\% \\
Wall unchanged & 17.2\% [11.4--27.2] & 27.7\% & 30 & -- & -- \\
Wall $= 1/q$ & 21.9\% [15.5--54.8] & 46.6\% & 29 & -- & -- \\
\midrule
Eq.~\ref{eq:hidden} with measured $s$ (not a prediction) & 4.3\% [3.1--5.8] & 6.4\% & 71 & -- & -- \\
\bottomrule
\end{tabular}}
\end{table}

\emph{A prediction from runs with no limit.} To know the cost before
running under a limit, we replay a run that had no limit. We cut the run
into slices of 0.2\,s, the interval at which we read the cgroup counters.
For each slice, the counters give the CPU seconds $d$ and the bytes read
$b$ of the agent's cgroup. A slice inside a model answer keeps its length
$\Delta t$. A local slice cannot finish faster than its CPU work allows
under a quota of $q$ cores. It also cannot finish faster than its reads
allow under a read limit of $r$ bytes per second. It therefore takes
$\max(\Delta t, d/q, b/r)$. The predicted wall time is the sum of the
slices. Nothing is fitted. A setting has 5 to 15 runs with no limit. We
replay each of them and take the median as the prediction for the
setting. This use of a recorded run is close to trace-driven noise
simulation in HPC~\cite{hoefler2010} and to trace-based capacity
planning~\cite{gmach2007}.

We designed the replay on the five CPU settings of \texttt{t3} in set~1.
We then tested it on the other 89 settings with a CPU quota or a read
limit: 44 on real tasks, 23 on other settings of \texttt{t3} (other
models, sets and machines), 11 on \texttt{t2} and 11 on the file read
\texttt{t8}.

\emph{Overall accuracy.} Figure~\ref{fig:hidden} shows every setting, and
Table~\ref{tab:predict} compares the replay with simpler predictions
that use the same inputs. The replay is within a median 6.9\% of the wall
time. It keeps the model slices as they are, so part of this accuracy
comes for free. We therefore also give the error on local time, the part
that the replay changes: a median of 7.1\%. The settings are not
independent, because settings with different limits on the same task
share the same runs with no limit. We therefore also compare predictions
per group of settings that share these runs (35 groups), with two-sided
sign tests.

\emph{The slices matter.} The same rule applied to the whole local part
at once is off by 11.8\% on local time. The replay is closer in 29 of the
35 groups ($p < 0.001$).

\emph{Against the simple estimate.} The simplest prediction is
Equation~\ref{eq:hidden} with $s = 1/q$. How the replay compares with it
depends on the limit.
\begin{itemize}
\item \emph{CPU quotas below one core (38 settings).} The simple estimate
  is as good or better: 7.6\% against 11.6\% on local time. The replay is
  closer in 13 of the 38 settings ($p = 0.07$). Under a tight quota,
  almost all of the agent's local work waits for the CPU, so
  $s \approx 1/q$.
\item \emph{Quotas of one core or more (37 settings).} Here $1/q$
  predicts no change. The replay is closer in 21 of the 28 settings where
  the two differ ($p = 0.01$).
\item \emph{Read limits (14 settings).} Here $1/q$ says nothing. The
  replay is within 4.7\% of local time, and the simple estimate is off by
  59\%. The replay is closer in all 14.
\end{itemize}
Over all settings and on wall time, the two do not differ significantly
once settings that share runs are grouped (the replay is closer in 19 of
35 groups, $p = 0.74$).

\emph{Where it holds.} The replay holds on each machine: its error on
wall time is 4.9\% on set~2 of M0, 7.6\% on M1 and 5.7\% on M2. It holds
on real tasks (7.0\%). It also explains why a CPU quota and a disk limit
together cost less than their sum: each slice waits only for the tighter
of the two.

\emph{Where it misses.} The replay misses the wall time by more than 10\%
in 28 settings. In 19 of them the model time changed by 2\,s or more,
which means the agent did something different. The other nine include
the lowest read limits and Django under 512\,MB, where memory reclaim
adds time that the replay cannot see. The long Sympy task of
Section~\ref{sec:rq3} shows this limit of the method more clearly. At 1
CPU with 512\,MB, the memory limit and not the quota sets the cost, and
both predictions miss it.

The median hides a tail, so Table~\ref{tab:predict} also gives the mean
and the 90th percentile. On wall time, the mean error of the replay is
11.6\%, and one setting in ten is off by 28\% or more. This tail is the
settings where the agent changed what it did. In the 58 settings where
the model time stayed within 2\,s, the median error is 5.2\%, the mean
6.4\% and the 90th percentile 11.5\%.

\emph{How many runs with no limit are needed?} The prediction above uses
all runs with no limit of a setting. We repeated it with each of these
runs alone. One run predicts the wall \emph{time} less well, because its
own model wait is one draw of a variable quantity. The typical run of a
setting is off by 9.7\% (median over settings), and the predictions from
different runs of one setting spread by 34\%. One run predicts the wall
\emph{slowdown} well, because its model wait is on both sides of the
ratio. The replayed time of a run divided by its own wall time is within
7.1\% of the measured slowdown for the typical run, and within 10.6\% for
the worst run of a setting. One run with no limit is therefore enough to
predict by how much a limit slows a task. Several runs are needed to
predict the time in seconds.

\textbf{Answer to RQ2.} Each limit has a breaking point and one of three
shapes: smooth (CPU, disk, threads), cliff (memory) or crash (process
count). For most limits the cost lands in the tools. For a disk limit on
the coding task it lands in start-up, and for the network only in model
time. The kernel trace shows that a CPU quota and a disk limit grow one
kind of waiting and leave the work unchanged, while a memory limit and a
busy neighbour also add CPU time. The model wait dilutes the cost by
Amdahl's law: when local time grows $s$ times, the user sees a slowdown
of $1 + f\,(s-1)$, with $f$ between 0.17 and 0.73 in our runs. For CPU
quotas below one core, this formula with $s = 1/q$ predicts the cost. For
disk limits, a replay of the runs with no limit is needed. It predicts
local time within a median 4.7\%, where the simple estimate is off by
more than half. Both hold as long as the agent does not change what it
does.

% ---------------------------------------------------------------------------
\section{RQ3: Do Real Platform Sizes Slow Down Real Tasks?}\label{sec:rq3}

\begin{table*}[t]
\caption{Real SWE-bench Verified tasks in Docker at platform sizes.
Slowdown against no limit: geometric mean of the per-task median ratios,
with 95\% bootstrap intervals. For the long tasks, the slowdown is the
geometric mean over the
three tasks that fit in 512\,MB; xarray is reported in the text.}
\label{tab:real}
\small
\begin{tabular}{@{}l ccc ccc@{}}
\toprule
& \multicolumn{3}{c}{6 short tasks} & \multicolumn{3}{c}{4 long tasks} \\
\cmidrule(lr){2-4}\cmidrule(l){5-7}
Size (platform) & Solved & Wall & Local & Solved & Wall & Local \\
\midrule
No limit & 30/30 & 1.00 & 1.00 & 11/20 & 1.00 & 1.00 \\
2 CPU, 8\,GB (CI runner) & 30/30 & 1.01 [0.86--1.14] & 0.98 [0.84--1.13] & -- & -- & -- \\
2 CPU, 512\,MB (sandbox) & 30/30 & 1.02 [0.90--1.14] & 0.99 [0.85--1.14] & -- & -- & -- \\
1 CPU, 512\,MB (Cloud Run) & 30/30 & 1.02 [0.90--1.38] & 1.06 [0.92--1.59] & 12/20 & 1.11 [0.58--1.39] & 1.25 [0.53--1.92] \\
0.25 CPU, 512\,MB (Fargate) & 30/30 & \textbf{1.74} [1.46--1.94] & \textbf{2.51} [2.00--2.87] & 11/20 & \textbf{2.12} [1.09--2.37] & \textbf{4.72} [1.96--5.51] \\
\bottomrule
\end{tabular}
\end{table*}

The small tasks of RQ2 show where each limit starts to hurt. We now ask
whether the sizes that real platforms offer hurt real tasks. We ran the
ten real tasks in Docker at the platform sizes of
Section~\ref{sec:design}. With no limit, a short task takes 26\,s and 4
model turns (medians), and a long task 64\,s and 9 turns. The short tasks
ran at all five sizes. The long tasks ran with no limit, at 1 CPU and at
0.25 CPU. Table~\ref{tab:real} gives the results.

\emph{1 CPU or more: no cost that we can detect.} At 1 CPU or more, every
interval in Table~\ref{tab:real} includes 1. The intervals are wide,
however: at 1 CPU, the one for wall time reaches 1.38. A cost below that
size is therefore possible. Means give a larger value at 1 CPU. With
per-task means instead of medians, the wall slowdown is 1.11 and not
1.02. The difference comes from one Requests run that took 123\,s,
against 20 to 24\,s for the other four. At the two sizes with 2 CPUs,
means and medians agree (0.96 to 1.02). The agent chooses different
tests in different runs, which adds variation. We therefore also
compared only the runs of the same test command. That command takes the
same time at 2 CPUs and 2 to 9\% longer at 1 CPU.

\emph{0.25 CPU: every task is slower.} At the Fargate minimum, every
short task is still solved. But its wall time is 1.74 times longer and
its local time 2.51 times longer (1.67 and 2.41 with means). A run takes
22\,s more. Most of this is in the tools (15\,s) and in start-up
(3.5\,s). Model time does not change (15.2 against 15.4\,s). The long
tasks spend more of their time on the machine, so more of the cost
reaches the user: local time grows 4.7 times and wall time 2.1 times.
With three tasks and five runs each, the interval for the long tasks is
wide (1.09 to 2.37 for wall time). We take from it only that these tasks
are slower. This matches the CPU breaking point of 0.75 core from RQ2:
1 CPU lies above it, and 0.25 CPU below.

\begin{table}[t]
\caption{Each short task at 0.25 CPU and 512\,MB, per machine: wall
slowdown against no limit on the same machine and set, and tasks solved.
Set~2 of M0 has 10 runs per setting, the others 5.}
\label{tab:pertask}
\small
\begin{tabular}{@{}lcccc@{}}
\toprule
Task & M0, set 1 & M0, set 2 & M1 & M2 \\
\midrule
Django 15930 & 1.70 (5/5) & 2.10 (10/10) & 2.38 (5/5) & 2.00 (5/5) \\
Sympy 24562 & 1.86 (5/5) & 1.85 (10/10) & 2.17 (5/5) & 2.11 (5/5) \\
Requests 5414 & 1.35 (5/5) & 1.57 (10/10) & 1.54 (5/5) & 1.57 (5/5) \\
Flask 5014 & 2.15 (5/5) & 1.55 (10/10) & -- & -- \\
pytest 10081 & 1.57 (5/5) & 1.68 (10/10) & -- & -- \\
xarray 6461 & 1.90 (5/5) & 2.09 (10/10) & -- & -- \\
\bottomrule
\end{tabular}
\end{table}

\emph{Per task and per machine.} Table~\ref{tab:pertask} gives each short
task at the Fargate size. The median wall time of every task is longer
on every machine, by 1.35 to 2.38 times, and every run is solved. With
five runs, the interval of a single task can still include 1 (Flask and
xarray in set~1, Requests on M1). At 1 CPU, no task on M1 or M2 is
clearly slower (0.85 to 1.05 times), except Django on M2 (1.37 times).

\emph{The two sets of M0 differ with no limit, and the trace shows
where.} Under the 0.25 CPU limit, local time is the same on every machine
and set (31 to 33\,s for Sympy). With no limit, local time is shorter on
set~2, M1 and M2 (5 to 8\,s) than on set~1 (8.5 to 13\,s). The kernel
trace locates the difference. In both sets, a run of a task reads the
same amount from disk in the same number of reads (for Django, 256 and
265\,MB in about 10,000 reads). But the threads wait for the disk for 3.5
to 11.8\,s of thread time in set~1, and for 1.2 to 3.1\,s in set~2. For
Django this is about 1.2\,ms against 0.3\,ms per read. Small reads were
thus slower on the disk of set~1, although its sequential read speed was
higher. Two other things also changed between the sets: the agent
version (Section~\ref{sec:design}), and the CPU time of the tools, which
is lower in set~2. We cannot say how much of the gap each change
explains. Under the 0.25 CPU limit, the wait for the CPU quota is much
larger than these differences, which is why the sets agree there. We
compare each setting only with runs from the same set.

\emph{Whether a task is solved depends on the task, not the size.} The
long Django and Sympy tasks are solved in all 30 of their runs. The long
pytest task is solved in 4 of 15 runs, about as often with and without
limits. The long xarray task (6992) is never solved, not even with no
limit. In its runs, all 902 tests that passed before still pass, but the
agent's fix does not give the result that the 12 target tests expect.
The reference fix does give it, in the same environment.

\emph{512\,MB is at the edge.} Five short tasks reach a peak memory of
345 to 453\,MB. The Django tasks reach 512\,MB, and the kernel takes back
memory without a kill. The test runs of the long xarray task need about
810\,MB. Under 512\,MB, every one of its runs has 4 to 9 memory kills.
The run takes 5.1 times longer at 1 CPU and 12 times longer at 0.25 CPU.
The agent never learns why (Section~\ref{sec:rq5}).

\emph{A long session.} The runs above last 26 to 64\,s. To see whether
the picture holds over many minutes, we ran a fifth long task in set~2:
Sympy issue 13878, rated at more than 4 hours for a human. It ran 5
times each with no limit, at 1 CPU with 512\,MB and at 0.25 CPU with
512\,MB. Every run solved the task.

With no limit, a run takes 300\,s (233 to 718\,s) and 29 model turns.
Its local share is $f = 0.73$. A task with a heavy test suite spends
most of its time on the machine, so the dilution is small.

At 1 CPU and 512\,MB, the run takes 2.3 times longer (local time 2.5
times). The short tasks showed no cost at this size, and the runs with
no limit rarely use more than one core. The counters show the cause. The
task's working set, the memory it uses actively, is 510 to 660\,MB. The
cgroup therefore hits the 512\,MB limit hundreds of times per run, and
the kernel takes memory back each time. No process is killed. But the
cgroup's CPU time grows from 234 to 568\,s. The agent makes 43 turns
instead of 29, after up to two of its tools time out.

At 0.25 CPU, the run takes 2.9 times longer (local time 3.4 times). Here
the quota sets the cost: local time is about the cgroup's CPU time
divided by the quota. The slowdown stays below $1/q = 4$ for one reason.
The agent's tool time-out, 2 minutes by default, stops long test runs 1
to 5 times per run. The agent then runs smaller tests, so the total CPU
work falls to 199\,s. Equation~\ref{eq:hidden} predicts the wall
slowdown at this size within 9\%, and the replay within 14\%. At 1 CPU,
both predict no cost, because the memory limit and not the quota sets
the cost. With five runs per size, and wall times that vary by a factor
of three with no limit, the intervals are wide (1 CPU: 0.95 to 3.7;
0.25 CPU: 1.2 to 5.2).

\textbf{Answer to RQ3.} Platform sizes of 1 CPU or more show no slowdown
that we can detect on the short tasks. The 0.25 CPU size slows every
task on every machine: 1.7 times for short tasks on average, about twice
for long tasks (with a wide interval), and 2.9 times in sessions of 14
to 24 minutes. The CPU breaking point predicts that this size is too
small, and Equation~\ref{eq:hidden} and the replay predict by how much.
The 512\,MB memory limit matters when a task needs more memory. On a
long task whose working set is near 512\,MB, the size with 1 CPU costs
2.3 times. On a task whose tests need about 810\,MB, the run takes 5 to
12 times longer. No size changes whether a task is solved. Every task
that is solved with no limit is solved at every size, and the two long
tasks that fail (xarray and pytest) fail as often with no limit.

% ---------------------------------------------------------------------------
\section{RQ4: How Many Agents Fit on One Machine?}\label{sec:rq4}

\begin{figure}[t]
\centering
\includegraphics[width=\columnwidth]{ICPE-2027/figures/fig_density.pdf}
\Description{Left: agents finished per minute per vCPU against agents per vCPU, for machines with 2, 4 and 8 vCPUs. The three curves fall on each other and level off at about 4 at four agents per vCPU, far below a dashed line at about 10, the estimate from one agent alone. Right: CPU seconds per agent rise from about 6 to 12 or 13 on all three machines.}
\caption{Many agents started at once on one machine, on three machine
sizes (\texttt{t3}; M0 from set~2). Both x-axes are agents per vCPU, so
the machines fall on one curve. Left: throughput per vCPU stops growing
at about four agents per vCPU. The dashed line is the CPU bound computed
from one agent alone. Right: the CPU time that each agent needs doubles as
the machine fills.}
\label{fig:density}
\end{figure}

RQ2 and RQ3 limit one agent at a time. A platform that runs many agents
shares one machine among them instead, so each agent's neighbours become
its limit. We started 1 to 32 Claude Code agents at the same moment on
one VM, with no limits. Each agent worked on its own copy of the coding
task \texttt{t3}. We call the agents that start together a \emph{group}.
All 189 agents solved their task. Crowding makes agents slow, not wrong.

We count throughput as the number of agents in a group divided by the
time from the first start to the last finish, in agents per minute. This
is the throughput of one batch. It is not the steady-state throughput of
a stream of arriving agents, which is what operational analysis bounds.
We first describe set~1 of M0, with three groups per level.
Figure~\ref{fig:density} then shows the same experiment on all three
machine sizes. A test with agents that arrive at random follows.

\emph{When the machine fills up.} Two agents run as fast as one (21 to
22\,s). At four agents, one per vCPU, local time grows by 41\% but wall
time only by 14\%. This is the dilution of Section~\ref{sec:rq2} again.
Throughput rises from 2.8 agents per minute with one agent to 16.3 with
sixteen, and then stops: 32 agents finish at 16.5 per minute. The extra
time is in the tools (4\,s alone, 71\,s at 32 agents) and in start-up
(1.4\,s alone, 15\,s at 32), where the agents wait for a CPU. Model time
stays the same up to 16 agents (16 to 17\,s) and is 20\,s at 32. The API
cost per agent stays at \$0.14 to \$0.15.

\emph{A capacity bound, and why the simple version fails.} Operational
analysis~\cite{denning1978} gives an upper bound on the throughput of $N$
agents: $\min(N/R_1,\ C/D)$. Here $R_1$ is the time of one agent alone,
$C$ the number of CPUs and $D$ the CPU seconds per agent. The first term
says that agents cannot finish faster than when each runs alone. The
second says that the CPUs cannot do more than $C$ seconds of work per
second. With the values of one agent alone ($R_1 = 21$\,s, $D = 5.9$\,s),
the CPU bound is about 40 agents per minute. This is 2.5 times what we
measured, because $D$ does not stay fixed. It is 8.3\,s with four agents
and 11.6 to 12.2\,s with sixteen or more. With $D$ measured at each
load, the measured throughput reaches 79 to 84\% of the bound.

A likely reason for the growth of $D$ is that each vCPU is a hardware
thread that shares a core and its caches with another. A full machine
then does less work per CPU second. The VMs give no access to hardware
counters, so we cannot measure this directly. The trace of the neighbour
experiment supports it. With four CPU-bound programs next to one agent,
the agent's threads run for 10.4\,s instead of 6.2\,s for the same task,
while they wait only 3.0\,s for a busy core (Section~\ref{sec:rq2}). The
extra time is spent running on a CPU, not waiting for one.

The same error appears in methods that plan capacity by adding up the
demand of each workload measured alone~\cite{gmach2007}. Sandbox
platforms place more work on a host than it can run at once. This works
when tenants do not peak together~\cite{agache2020}. Identical agents on
similar tasks tend to peak together.

\emph{CPU or memory?} Each extra agent adds about 150\,MB. On this 15\,GB
machine, memory would last to about 100 agents, but the CPU is full at
16. A study of agents in larger sandboxes estimated that memory limits
how many agents fit, because an agent's peak memory there can reach 2 to
4\,GB~\cite{agentcgroup}. It computed this from single-agent peaks and
did not run the agents together. Both results follow one rule. Divide the
machine's memory by the number of agents at which the CPU is full. If one
agent's peak memory is above this value, memory runs out first. Here the
value is $15\,\text{GB}/16 \approx 0.9$\,GB. Our real tasks peak at 345
to 811\,MB with no limit (medians), and one run of the long xarray task
reached 972\,MB.

\emph{Other machines and another task.} We repeated the experiment on M1
(1 to 16 agents, 3 groups per level), on M2 (1 to 64 agents, 3 groups)
and on set~2 of M0 (1 to 32 agents, 5 groups; 3 at two agents).
Figure~\ref{fig:density} shows the three machines. Throughput stops
growing at about four agents per vCPU on every machine:
\begin{itemize}
\item M1 (2 vCPUs): 7.7 agents per minute at 8 agents, and 7.5 at 16.
\item M0, set~2 (4 vCPUs): 14.5 per minute at 16 agents. Going to 32
  agents adds 6\%.
\item M2 (8 vCPUs): 32.4 per minute at 32 agents. Going to 64 agents
  adds 5\%.
\end{itemize}
On every machine, the CPU seconds per agent grow from about 6 alone to 12
or 13 on a full machine. The CPU bound computed from one agent alone is
2.3 to 2.5 times the measured throughput. Each extra agent adds about
150\,MB on every machine. The slowest agents stay close to the median:
once a machine is full, the 95th percentile of the time per agent is 1.04
to 1.14 times the median.

A second, lighter task (Requests in Docker, 4.5\,s of CPU alone) still
gains throughput at 16 agents on M0 (19.5 per minute). The point where a
machine is full therefore depends on the task. Over all machines, sets
and both tasks, all 1,067 agents solved their task.

\emph{Agents that arrive at random.} Agents that start at the same
moment are a worst case, because their busy moments coincide. We
therefore repeated the experiment on set~2 of M0 with 4 to 32 agents
whose start times are spread at random over 60\,s. There are 5 groups
per level and 300 agents, and all solved their task. The result does not
change.
\begin{itemize}
\item \emph{4 to 16 arrivals per minute.} The arrivals stay below the
  machine's rate, and the agents hardly overlap (1 to 4.5 alive at once
  on average). Each runs as if alone: local time is 4.6 to 6.4\,s,
  against 4.7 to 5.7\,s for one agent.
\item \emph{32 arrivals per minute.} The agents pile up. On average 14.5
  are alive at once, and 24 at the peak. Each agent's local time is
  42\,s, the same as in the simultaneous groups of 16 (40\,s). The CPU
  time per agent is 11.6\,s, against 6.2\,s alone. The throughput settles
  at 16.9 agents per minute, 10\% above the 15.4 of the simultaneous
  groups of 32.
\end{itemize}
The point where the machine is full, and the extra CPU time per agent,
therefore do not depend on the agents starting together.

\textbf{Answer to RQ4.} Up to one agent per vCPU costs the user little
(14\% more wall time on 4 vCPUs). Throughput stops growing at about four
agents per vCPU, on 2, 4 and 8 vCPUs, for this task, whether the agents
start together or arrive at random. CPU is the limit, not memory. The
CPU time that each agent needs doubles when the machine is full. A
capacity estimate from one agent alone is therefore 2.3 to 2.5 times too
high on every machine we tried.

% ---------------------------------------------------------------------------
\section{RQ5: Does the Agent Change Its Behaviour?}\label{sec:rq5}

RQ1 to RQ4 treat the agent as a fixed workload that the system slows
down. But an agent reads what the system shows it, and can act on it. We
looked for three such changes: extra work caused by a slow result,
reactions to system errors, and lighter tests on a slow machine.

\emph{A printed slow time makes a small model work more.} The coding task
\texttt{t3} prints its own run time. Task \texttt{t9} is the same task
with that line removed. At 0.25 core, we ran each task 15 times with
Haiku. When the time is printed, Haiku talks about speed in 15 of 15
runs. It keeps trying to make the code faster after the tests already
pass. It makes 5 tool calls after the tests first pass, against 0 when
no time is printed. It also takes 11 model turns against 8, and 73\,s
against 35\,s of wall time ($p \leq 0.04$ after a Holm correction over
the 20 tests per model). With no limit, the printed time causes no extra
work. At 0.5 core the effect is not significant. Opus does no extra work
in any setting.

The effect depends on the task. On the second coding task \texttt{t10},
the test prints a smaller time. Haiku mentions speed in 5 of 5 runs at
0.25 core but does not act on it. It makes 0 tool calls after the tests
pass, and 9 turns against 7 ($p = 0.15$).

This extra work explains two results of RQ2. At 0.25 core, Haiku uses
2.9 times more CPU time than with no limit (16.8 against 5.8\,s), and it
takes extra model turns. A slow machine thus costs a small model twice:
once directly, and again through the extra work that the slow number
causes. The extra work also costs money. In set~2, a Haiku run at 0.25
core writes 8,600 output tokens when the time is printed and 4,100 when
it is not, and its API cost is 44\% higher (medians). For Fable the cost
is 20\% higher.

\emph{Controls.} We then tested the cause in set~2, with the three
variants of task \texttt{t11} (Section~\ref{sec:design}). Each
comparison has 10 or 15 runs per side, with a Holm correction over 18
tests. All results are for Haiku at 0.25 core
unless we say otherwise.
\begin{itemize}
\item \emph{The effect repeats in set~2.} Haiku makes 4 tool calls after
  the tests first pass when the time is printed, and 0 when it is not. It
  takes 79\,s against 30\,s ($p_\text{Holm} = 0.02$ for the tool calls).
\item \emph{A fixed fast time removes it.} When the test prints a fixed
  fast time (2.5\,s) while the machine is just as slow, the extra work
  disappears (0 calls, 33\,s).
\item \emph{One sentence in the prompt removes it.} The sentence says
  that the printed time is for information only and is not a goal (0
  calls, 31\,s).
\item \emph{A fixed slow time alone does not cause it.} A fixed slow
  time (12\,s) on a machine with no limit gives 0 calls and 17\,s.
\item \emph{Fable shows the effect more weakly.} It talks about speed in
  13 of 15 runs, makes 1 extra tool call against 0, and takes 64\,s
  against 37\,s ($p_\text{Holm} < 0.001$).
\end{itemize}
The trigger is therefore neither the slow machine alone nor the slow
number alone. One explanation, which we did not test, is that the agent
chases a number that changes when it edits the code. A fixed number
gives it nothing to chase. Figure~\ref{fig:feedback} shows every run of
these tests.

\begin{figure}[t]
\centering
\includegraphics[width=\columnwidth]{ICPE-2027/figures/fig_feedback.pdf}
\Description{Dot plots of six groups of Haiku runs. Left: tool calls after the tests first pass. Right: wall time. Only the group with the real time printed at a quarter of a core has extra tool calls (median 4) and a long wall time (median 79 seconds). The other five groups have a median of 0 extra calls and 17 to 33 seconds.}
\caption{The printed run time and its controls (Haiku, set~2, 10 or 15
runs per group). Each dot is one run; the bar is the median. Only a real
slow time, printed on a slow machine, causes extra work.}
\label{fig:feedback}
\end{figure}

\emph{The agent reports the errors it can see, and only those.} When a
memory limit kills the tool, the agent sees the kill (exit code 137,
which means the process was killed). It names the cause in its final
message in 23 of 23 runs. It does not retry in a loop: it makes 0 to 2
tool calls after the error. This agrees with work showing that explicit
reports of what happened reduce failures of state
awareness~\cite{song2025}. Three cases leave the user with no
explanation.

\emph{Case 1: the kernel kills the agent instead of the tool.} The
kernel's out-of-memory killer picks one process in the cgroup, and the
agent itself is a candidate. With the 512\,MB tool of \texttt{t7}, the
agent was killed in 2 of 25 runs of set~1 (at 576 and 512\,MB) and in 5
of 20 runs at 640\,MB on set~2, M1 and M2. With the 700\,MB test of
\texttt{t12} under a 768\,MB limit, it was killed in 9 of 20 runs (Haiku
6 of 10, Opus 3 of 10). It was never killed in the 20 runs with no
limit. The run then ends without a message. In \texttt{t12}, the agent's
code was already correct in every one of these 9 runs.

The kernel's own report, which the trace records, shows why. At every
kill it lists the processes in the cgroup with their memory. In all 60
kills in \texttt{t12}, the kernel picked the largest process. The first
kill of a run always hits the test (about 615\,MB). The agent holds
about 150\,MB. In the 9 runs where the agent died, the cgroup reached
its limit again within seconds after the test was killed. At that moment
no large process was left, so the largest one was the agent. We are not
sure why the cgroup was still at its limit. One possible reason is
memory that the kernel had not yet taken back from the dead test. A test
that needs most of a memory limit is not rare in real projects: the long
xarray task needs about 810\,MB and ran under a 512\,MB platform size
(Section~\ref{sec:rq3}).

The usual ways to protect a process from the out-of-memory killer are
known. One is a separate cgroup for the tools. Another is
\texttt{oom\_score\_adj}, a per-process value that tells the killer
which process to prefer. A third is \texttt{memory.oom.group}, which
makes the kernel kill a whole cgroup together. AgentCgroup proposes one
cgroup per tool call~\cite{agentcgroup}. We tested two of these fixes on
the same \texttt{t12} setting, with 10 runs per model and fix.
\begin{itemize}
\item \emph{Fix 1.} Every shell that the agent starts for a tool raises
  its own \texttt{oom\_score\_adj} to the maximum, so the kernel prefers
  to kill tools.
\item \emph{Fix 2.} Every tool shell moves into a child cgroup with its
  own \texttt{memory.max} of 448\,MB, so a kill in that cgroup can only
  hit tools.
\end{itemize}
With either fix, the kernel killed the agent in 0 of 20 runs (95\%
Wilson interval 0 to 16\%), against 9 of 20 with no fix (one-sided
Fisher exact test, $p < 0.001$ for each fix). The kernel still killed
the test, and every run solved the task. The agent's final report named
the kill in 39 of 40 runs. The first fix made runs slower (27 to 29\,s
against 20 to 24\,s, medians). The second did not (19 to 21\,s).

\emph{Case 2: a process limit crashes the agent.} In set~1, between 17
and 20 processes, the trace shows 1 to 6 failed attempts to start a
process per run. The shell retries these attempts, so the agent sees an
error in only 3 of 20 runs. But the agent program also needs threads of
its own. When it cannot start one, it aborts. This happened in 3 of 15
runs at 17 to 19 processes, and in 9 of 10 runs at 16 and 12.

\emph{Case 3: memory kills inside a test run are silent.} In the xarray
runs under 512\,MB (Section~\ref{sec:rq3}), the kernel kills test
workers. The tool output simply stops partway (for example at
``34\%''), with no error message and no exit code. In 6 of 10 runs, the
agent saw no error at all, and in none did its final answer mention
memory. Instead, it ran the tests again (10 or 11 turns against 4 with
no limit). SWE-bench counts a test that reports no result as
failed~\cite{jimenez2024}. Such a kill can therefore lower the measured
solve rate even when a fix is correct.

The small test of \texttt{t12} shows the other outcome. There the output
stops at 50\%. In 16 of 20 runs the agent ran the test again without
\texttt{tee} and saw exit status 137. In all 11 runs where the agent
itself was not killed, it reported that the test was killed and did not
pass.

\emph{The agent does not adapt its tests to a slow machine.} On the short
real tasks at 0.25 CPU, Opus runs the same number of tests as with no
limit. In the 60 runs of set~1 on these tasks, with no limit and at 0.25
CPU, the agent never mentions slowness, limits or time-outs. When no time
is printed, the agent does not notice that the machine is slow.

In long sessions, the agent's own tool time-out becomes a signal. On the
long Sympy task (Section~\ref{sec:rq3}), the Bash tool stops a command
after 2 minutes by default. This happened 0 to 2 times per run with no
limit and 1 to 5 times per run under the two limits. The agent then ran
smaller sets of tests, made more turns (43 to 56 against 29), and still
solved the task in every run.

\textbf{Answer to RQ5.} Yes, in two ways. First, a small model reacts to
a slow printed time by doing extra work, which makes the run 2.1 to 2.6
times longer. Opus does not do this, Fable does it more weakly, and one
sentence in the prompt prevents it. Second, the system can hide a
failure. A memory kill that hits the agent leaves the user with no
message, even when the work is done. So does a process limit that
crashes the agent. A kill inside a large test run can go unnoticed. A
kill that the agent can see, on a tool or a small test, is explained
correctly almost every time (73 of 74 runs). Keeping the out-of-memory killer away from the
agent removed the kills of the agent in our tests. Two fixes did this: a
child cgroup for the tools, and a higher \texttt{oom\_score\_adj} for
them.

% ---------------------------------------------------------------------------
\section{Results Across Machines}\label{sec:cross}

\begin{table*}[t]
\caption{The same experiments on three machine sizes. Breaking points are
on local time (Section~\ref{sec:design}). M0 values are from set~1 unless
a range covers both sets. The last column says what sets each number.}
\label{tab:cross}
\small
\begin{tabular}{@{}lllll@{}}
\toprule
Result & M1 (2 vCPU, Intel) & M0 (4 vCPU, Intel) & M2 (8 vCPU, AMD) & Set by \\
\midrule
CPU quota: breaking point & 0.75 core & 0.75 core & 0.75 core & the task \\
CPU quota 0.25 core: local (wall) & 4.8$\times$ (2.6$\times$) & 3.9--4.0$\times$ (1.8--2.0$\times$) & 4.0$\times$ (1.9$\times$) & the task \\
Threads: breaking point & 2 per vCPU & 2 per vCPU & 2 per vCPU & vCPUs \\
Soft memory limit: cliff & 1152\,MB & 1152\,MB & 1152\,MB & file size (1\,GB) \\
512\,MB tool: killed at & 640\,MB & 640\,MB & 640\,MB & tool and agent memory \\
Disk read limit: breaking point & 150\,MB/s & 150\,MB/s & 50--150\,MB/s & disk speed \\
Process limit: agent crashes at & 12 & 16 & 24 & vCPUs \\
Processes and threads kept alive, no limit & 12--13 & 16--17 & 24--26 & vCPUs \\
Many agents: throughput stops at & 8 agents & 16 agents & 32 agents & vCPUs and task \\
\bottomrule
\end{tabular}
\end{table*}

The results so far come mostly from M0. To see which of them depend on
the machine, we repeated a core set of experiments on two other VMs
(Section~\ref{sec:design}). M1 has half of M0's CPUs and memory. M2 has
twice its CPUs, from another vendor. We also compare M0 with VM0, a
first VM of the same type as M0 with an older kernel and agent version.
Table~\ref{tab:cross} gives the result.

\emph{What stays the same.} The mechanisms hold on every machine.
\begin{itemize}
\item \emph{CPU quota.} The breaking point is 0.75 core on 2, 4 and 8
  vCPUs. The quota limits the agent's cgroup, not the machine. What
  matters is therefore how much CPU the task needs in its busiest
  moments. The effect on wall time is also large on every machine. At
  0.25 core, Cliff's $\delta$ against no limit is 1.0 on all four sets
  of runs, which means that every limited run took longer than every run
  with no limit. At 0.5 core it is 0.84 to 1.0. At 0.75 core it is 0.2
  to 0.8, so the wall effect there is small or unclear even where local
  time grows.
\item \emph{Memory.} The memory cliffs are set by the task. The 1\,GB
  file of \texttt{t6} no longer fits in memory at a soft limit of
  1152\,MB, and the 512\,MB tool of \texttt{t7} is killed at 640\,MB, on
  every machine.
\item \emph{Extra threads.} They start to cost at two workers per vCPU
  everywhere. With $W$ workers on $C$ cores, the share of time spent
  waiting for a CPU follows $(W-C)/W$. With 8 workers on 4 vCPUs it is
  0.48 against the expected 0.50, and with 16 workers 0.73 against 0.75
  (VM0 and M0).
\end{itemize}

\emph{What scales with the machine.} Two numbers follow the number of
vCPUs. The first is the process limit at which the agent crashes. It is
12 on M1, 16 on M0 and 24 on M2. At 14, 17 and 28, every run solves the
task.

The cause is the agent itself, and the kernel trace names it. The trace
gives the name of every thread in the agent's cgroup. The agent's
runtime (Bun with the JavaScriptCore engine) starts one worker thread
per vCPU (\texttt{Bun Pool}: 2, 4 and 8 threads). It also starts one
garbage collection helper per vCPU minus one (\texttt{HeapHelper}: 1, 3
and 7). About eight other threads are always there. Every thread counts
against \texttt{pids.max}.

Two counts describe what the agent needs. The first is the number of
processes and threads that the agent keeps alive, which we read 5 times
per second. With no limit, it is 12 to 13 on M1, 16 to 17 on M0 and 24
to 26 on M2, or about 8 plus 2 per vCPU. The second is the kernel's own
peak counter (\texttt{pids.peak}). It is higher (15, 21 and 33), because
of short-lived threads at start-up. The crash edge follows the first
count: limits of 14, 17 and 28 pass although they are below the second.
The agent crashes (exit code 134) almost always when the limit is below
the number of processes and threads it keeps alive. The exceptions are a
limit of 10 on M1, where 3 of 5 runs crashed, and the rare aborts above
the edge (Section~\ref{sec:rq2}).

The count of 8 plus 2 per vCPU describes three machine sizes and one
agent, so it is a summary and not a law. What carries over is the rule
behind it: the agent crashes when the limit is below the number of
processes and threads it keeps alive, and this number grows with the
vCPUs. We found this rule in the first runs on M1 and M2. We then tested
it with follow-up runs below 16 on M1 and above 24 on M2.

The second number that follows the vCPUs is the point where many agents
fill the machine: about four agents per vCPU (Section~\ref{sec:rq4}).

\emph{What depends on details of the machine.} The disk breaking point
follows the disk's speed. A 200\,MB/s limit had no effect on VM0, whose
disk reads 205\,MB/s. On M0, whose disk read 287\,MB/s in set~1, the
same limit costs 1.1 times. The long scheduling delays depend on the
kernel. With 16 workers, a run on M0 has 421 waits for a CPU of 10\,ms
or more, against 133 on VM0.

\textbf{Summary.} The mechanisms carry over across two CPU vendors and
three machine sizes. Thresholds that the task sets stay the same.
Thresholds that the agent's runtime or the machine sets move. One run
with no limit on the target machine measures how many processes and
threads the agent needs. The number of agents that fill a machine must
be measured with the machine full, because an estimate from one agent is
too high.

% ---------------------------------------------------------------------------
\section{Implications}\label{sec:implications}

Our results lead to concrete advice for three groups: those who measure
agents, those who build the platforms that run them, and those who build
the agents.

\emph{For anyone who measures agents.} Report local time next to wall
time, or run many more repetitions. With five runs, a change in wall
time below about 20\% is lost in the variation of the model's answer
time. A cost of 50\% on the machine reaches the user as only 10 to 15\%
on the small tasks ($f$ from 0.20 to 0.30), and as 9 to 37\% on the real
tasks ($f$ from 0.17 to 0.73). Benchmarks and A/B tests that compare
only wall time can therefore call such a setting free.
Equation~\ref{eq:hidden} converts one slowdown into the other, and $f$
needs only a few runs with no limit.

The share $f$ also depends on the agent. Some agent designs spend most
of a run on local work such as tests~\cite{fan2025}. This raises $f$, so
the same limit costs such agents more. Methods that shorten the model's
context cut the tokens of each run~\cite{lindenbauer2025}. If they also
shorten the model time, they raise $f$ in the same way. We did not
measure this.

\emph{For platform builders.} The numbers below hold for the machines
and tasks we measured. Each comes with a way to find the number for
another setup.
\begin{enumerate}
\item \emph{Give each agent at least one CPU.} In our runs, the CPU
  breaking point was 0.75 core on every machine. No size at or above 1
  CPU showed a cost that we could detect, although our intervals allow
  up to about 38\% at 1 CPU. A size of 0.25 CPU made short real tasks
  1.35 to 2.4 times longer. For another task, measure $f$ in a few runs
  with no limit and use Equation~\ref{eq:hidden} with $s = 1/q$. For a
  disk limit, replay one of these runs instead (Section~\ref{sec:rq2}).
  Autoscalers that set limits from past use~\cite{rzadca2020} can use
  such an estimate to see what a limit costs in time.
\item \emph{Size memory from the project's test suite, not from the
  agent.} The agent needs about 300\,MB, but one test suite needed
  810\,MB. Under a 512\,MB limit, its runs took 5 to 12 times longer.
\item \emph{Set the process limit above the agent's peak number of
  processes and threads on the target machine, plus room for the tools.}
  One run with no limit measures this peak (\texttt{pids.peak}), and a
  limit above it is safe. In our runs the peak was 15 to 33 on 2 to 8
  vCPUs. A limit that is safe on 2 vCPUs (16) therefore crashes the
  agent on 8. Common container runtimes set no process limit by default,
  so this matters where a platform sets one on purpose.
\item \emph{Keep the agent out of reach of the out-of-memory killer.} In
  our runs, the agent and its tools shared one cgroup. The kernel killed
  the agent itself in 7 of 45 runs with the 512\,MB tool at 640\,MB or
  less, and in 9 of 20 runs with a 700\,MB test at 768\,MB. Two cheap
  fixes each removed this in our tests (0 of 20 runs each). The first
  runs the tools in a child cgroup with their own \texttt{memory.max}.
  The second raises the \texttt{oom\_score\_adj} of the tools. We prefer
  the child cgroup. It cost no time in our runs, and it also gives the
  tools a memory limit of their own.
\item \emph{To pack agents on a machine, measure the CPU time per agent
  with the machine full.} With identical agents started together on our
  three machine sizes, this time doubled, and the machine was full at
  about four agents per vCPU. The value from one agent alone overstated
  the capacity 2.3 to 2.5 times.
\end{enumerate}

\emph{For agent builders.}
\begin{enumerate}
\item \emph{Tell the model when a tool was killed.} A child process that
  ends by \texttt{SIGKILL}, the signal that the kernel uses to kill a
  process, is something the agent program can see but the model cannot.
  The same holds for output that stops without an exit status. Tools
  that report state changes explicitly reduce failures of state
  awareness~\cite{song2025}.
\item \emph{Handle a failed thread start instead of aborting.} The agent
  aborted almost every time the process limit was below the number of
  processes and threads that it keeps alive.
\item \emph{Expect small models to act on numbers in tool output.} A
  printed run time made Haiku's runs 2.1 to 2.6 times longer, because of
  work that nobody asked for. In our runs, one sentence in the prompt
  saying that run time is not a goal removed this.
\end{enumerate}

% ---------------------------------------------------------------------------
\section{Threats to Validity}

\emph{Construct validity.} Local time is wall time minus the time during
which a model turn is open, taken from the agent's event stream. Work
that runs while the model still answers counts as model time, so local
time is a lower bound. The trace shows that this work is 0.25\% of local
time in the median run. Moving it to local time leaves the error of
Equation~\ref{eq:hidden} unchanged (Section~\ref{sec:design}). Tracing
costs about 7\% CPU time on the coding task. This cost is present on
both sides of every ratio. Tracing did not change wall time with no
limit, at 0.25 core, under a memory limit, or with 8 agents.

\emph{Internal validity.} The speed of the model service changes over
the day. We therefore run the steps of an experiment in random order
with a fixed seed. We also report local time, which does not contain the
model wait. The runs with no limit show why. In no setting did local
time change with the start time of the run. Wall time did in some
settings (Spearman $\rho$ up to 0.85), although no trend remains after a
Holm correction over the 18 tests. We compare runs only within one
machine and set, because local time with no limit differs between sets
for some real tasks (Section~\ref{sec:rq3}).

The replay was designed on five settings and tested on 89 others, so its
error is not tuned to the test data. It treats the CPU quota as a smooth
rate. The kernel enforces the quota per period: a cgroup that uses its
quota early in a 100\,ms period stops until the next
one~\cite{turner2010,cgroupv2}. Part of the replay's error may come from
this.

The long xarray task was never solved, but not because of our setup. As
for every task, its target tests fail at the starting version and pass
with the reference fix in the same environment.

\emph{External validity: machines.} We used four VMs of three types,
with two CPU vendors and 2 to 8 vCPUs, in one cloud. The mechanisms
agree on all of them, but some numbers scale with the vCPUs or depend on
the disk and kernel (Section~\ref{sec:cross}). We give these numbers as
rules to measure, not as fixed values. M2 differs from M0 in both the
number of vCPUs and the CPU vendor, so the two cannot be separated on it
alone. Two observations speak against an effect of the vendor. For the
process limit, the thread names in the trace tie the crash edge to
thread pools sized by the number of vCPUs. For density, the Intel
machines M1 and M0 and the AMD machine M2 fall on the same curve per
vCPU. We did not test whether hardware threads (SMT) cause the extra CPU
time per agent on a full machine.

We ran the agent in cgroups and Docker. Sandboxes such as Firecracker
share CPU through the same host cgroups~\cite{agache2020}, so the CPU
results should carry over. Memory and disk may not. In gVisor,
allocating memory is several times slower than on the
host~\cite{anjali2020}. Agents on one machine also share kernel locks,
for example in the page allocator, which cgroups do not
separate~\cite{anjali2025}. We did not measure lock waits, so they may
be part of the extra CPU time per agent on a full machine.

\emph{External validity: workload.} The workload is narrow in four ways.
\begin{itemize}
\item \emph{One agent.} All runs use one agent (Claude Code), which is
  closed source, and most use one model (Opus). Agent versions change
  every few weeks, and two versions appear in our runs. We expect the
  mechanisms (quota, reclaim, kills) to carry over to other agents,
  because the kernel enforces them, not the agent. The local share $f$
  and the process-limit rule depend on the agent and must be measured
  again. Agents that run more tests have a larger $f$~\cite{fan2025}, so
  the same limit would cost them more.
\item \emph{Few and short real tasks.} The real tasks are eleven Python
  issues from one benchmark, and most runs are short: 26\,s and 4 model
  turns for a short task, 64\,s and 9 turns for a long one. The
  benchmark is public, so the models may have seen these issues in
  training. That would make runs shorter and $f$ smaller than on unseen
  tasks. One task gives sessions of 4 to 24 minutes and 29 to 56 turns
  (Section~\ref{sec:rq3}). There the agent's tool time-out changed the
  run under limits, and the memory limit set the cost at 1 CPU. We have
  this for one task only.
\item \emph{Identical agents in the density experiment.} All agents work
  on the same task. Agents that start together and agents that arrive at
  random within a minute fill the machine at the same point. But the
  lighter Requests task fills the machine later, and mixed tasks may fit
  in larger numbers.
\item \emph{Small tasks for the behaviour results.} The extra work after
  a printed run time appeared on one of two small coding tasks. We did
  not test it on real tasks, and its cause is not established. We tested
  the two fixes for the out-of-memory killer on one task only.
\end{itemize}

\emph{Conclusion validity.} Set~1 has five runs per setting. Set~2 has
10 to 15 for the CPU quota, the process limit and the real tasks, and M1
and M2 have 5 to 10. We name a breaking point only where the whole
interval lies above 1.10. At 5\% or 20\% instead, the local breaking
points of the CPU quota, the soft memory limit and the process limit
stay the same on every machine. Of all 42 breaking points, 24 stay the
same, 15 move by one step and 3 by more.

With five runs, a non-parametric 95\% interval of the median cannot be
formed~\cite{hoefler2015}, so the bootstrap interval that we report may
cover less than 95\%. We checked this by drawing five runs from our
settings with 10 to 15 runs. The interval held the value from all runs
in 99\% of the draws. The rule missed some breaking points near the
threshold but never found a false one (Section~\ref{sec:design}).

The many-agent experiment has three groups per level in set~1 of M0, on
M1 and on M2, and five in set~2 of M0. The groups with random arrivals
were run to test the limit found in the first groups. Behaviour tests
use 10 to 15 runs per side and a Holm correction. We mark which results
are confirmatory and which are exploratory (Section~\ref{sec:design}).

% ---------------------------------------------------------------------------
\section{Conclusion}

An AI coding agent waits for its model most of the time. This dilutes
the machine's cost but does not remove it. We measured this in 2,199
runs on four VMs, with ten kinds of limits and eleven real tasks. Only a
fraction $f$ of a cost on the machine reaches the user, as Amdahl's law
predicts. Large costs stay visible, and moderate costs fall below the
variation of the model's answer time.

Kernel traces show that a CPU quota or a disk limit adds waiting and
leaves the agent's work unchanged. The cost can therefore be predicted
from runs with no limit. For CPU quotas below one core, the simple
formula with $s = 1/q$ predicts it. For disk limits, a replay of the
resource use of these runs in 0.2-second slices predicts local time
within about 5\%, on settings, machines and tasks that it was not
designed on.

Each limit has a breaking point and one of three shapes. The mechanisms
hold on every machine. Two numbers scale with the vCPUs. The first is
the process limit that crashes the agent, which one run with no limit
measures. The second is the number of agents that fill a machine, which
must be measured on a full machine.

Limits also change what the agent does. A small model does extra work
when it sees a slow time. The kernel can kill the agent without anyone
being told, and a separate cgroup for the tools prevented this in our
tests. Our dataset and collector let others repeat these measurements on
new agents and machines.

\section*{Data Availability}
Traces, collector and analysis scripts: \todo{anonymous link}.

\begin{thebibliography}{99}

\bibitem{petrini2003}
F.~Petrini, D.~J. Kerbyson, and S.~Pakin, ``The case of the missing
supercomputer performance: Achieving optimal performance on the 8,192
processors of ASCI~Q,'' in \emph{Proc.\ SC}, 2003.

\bibitem{tsafrir2005}
D.~Tsafrir, Y.~Etsion, D.~G. Feitelson, and S.~Kirkpatrick, ``System noise,
OS clock ticks, and fine-grained parallel applications,'' in \emph{Proc.\ ICS}, 2005, pp.\ 303--312.

\bibitem{ferreira2008}
K.~B. Ferreira, P.~Bridges, and R.~Brightwell, ``Characterizing application
sensitivity to OS interference using kernel-level noise injection,'' in \emph{Proc.\ SC}, 2008.

\bibitem{hoefler2010}
T.~Hoefler, T.~Schneider, and A.~Lumsdaine, ``Characterizing the influence
of system noise on large-scale applications by simulation,'' in \emph{Proc.\ SC}, 2010.

\bibitem{larres2013}
J.~Larres, A.~Potanin, and Y.~Hirose, ``A study of performance variations in
the Mozilla Firefox web browser,'' in \emph{Proc.\ ACSC}, 2013, pp.\ 3--12.

\bibitem{felter2015}
W.~Felter, A.~Ferreira, R.~Rajamony, and J.~Rubio, ``An updated performance
comparison of virtual machines and Linux containers,'' in \emph{Proc.\ IEEE ISPASS}, 2015, pp.\ 171--172.

\bibitem{maricq2018}
A.~Maricq, D.~Duplyakin, I.~Jimenez, C.~Maltzahn, R.~Stutsman, and R.~Ricci,
``Taming performance variability,'' in \emph{Proc.\ USENIX OSDI}, 2018.

\bibitem{mytkowicz2009}
T.~Mytkowicz, A.~Diwan, M.~Hauswirth, and P.~F. Sweeney, ``Producing wrong
data without doing anything obviously wrong!'' in \emph{Proc.\ ASPLOS}, 2009, pp.\ 265--276.

\bibitem{hoefler2015}
T.~Hoefler and R.~Belli, ``Scientific benchmarking of parallel computing
systems: Twelve ways to tell the masses when reporting performance
results,'' in \emph{Proc.\ SC}, 2015.

\bibitem{desnoyers2006}
M.~Desnoyers and M.~R. Dagenais, ``The LTTng tracer: A low impact
performance and behavior monitor for GNU/Linux,'' in \emph{Proc.\ Ottawa Linux Symp.}, vol.~1, 2006, pp.\ 209--224.

\bibitem{giraldeau2016}
F.~Giraldeau and M.~Dagenais, ``Wait analysis of distributed systems using
kernel tracing,'' \emph{IEEE TPDS}, vol.~27, no.~8, pp.\ 2450--2461, 2016.

\bibitem{ezzati2018}
N.~Ezzati-Jivan, G.~Bastien, and M.~Dagenais, ``High latency cause
detection using multilevel dynamic analysis,'' in \emph{Proc.\ IEEE SysCon}, 2018, pp.\ 1--8.

\bibitem{jimenez2024}
C.~E. Jimenez \emph{et al.}, ``SWE-bench: Can language models resolve real-world GitHub
issues?'' in \emph{Proc.\ ICLR}, 2024.

\bibitem{yang2024}
J.~Yang \emph{et al.}, ``SWE-agent: Agent-computer interfaces enable automated software
engineering,'' in \emph{Proc.\ NeurIPS}, vol.~37, 2024.

\bibitem{claudecode}
Anthropic, ``Claude Code documentation,''
\url{https://code.claude.com/docs/en/overview}, 2026.

\bibitem{codexcli}
OpenAI, ``Codex CLI,''
\url{https://github.com/openai/codex}, 2026.

\bibitem{cgroupv2}
Linux kernel documentation, ``Control Group v2,''
\url{https://docs.kernel.org/admin-guide/cgroup-v2.html}, 2026.

\bibitem{turner2010}
P.~Turner, B.~B. Rao, and N.~Rao, ``CPU bandwidth control for CFS,'' in \emph{Proc.\ Ottawa Linux Symp.}, 2010, pp.\ 245--254.

\bibitem{harchol2013}
M.~Harchol-Balter, \emph{Performance Modeling and Design of Computer
Systems}.\ Cambridge Univ.\ Press, 2013.

\bibitem{denning1968}
P.~J. Denning, ``The working set model for program behavior,''
\emph{Commun.\ ACM}, vol.~11, no.~5, pp.\ 323--333, 1968.

\bibitem{dean2013}
J.~Dean and L.~A. Barroso, ``The tail at scale,'' \emph{Commun.\ ACM}, vol.~56, no.~2, pp.\ 74--80, 2013.

% TODO: check every URL below and the arXiv id of earlier2026 before submission.
\bibitem{agentcgroup}
Y.~Zheng, J.~Fan, Q.~Fu, Y.~Yang, W.~Zhang, and A.~Quinn,
``AgentCgroup: Understanding and controlling OS resources of AI agents,''
arXiv:2602.09345, 2026.

\bibitem{earlier2026}
\textit{Cited in the third person; authors and arXiv id to be filled in},
``OS noise for agentic AI: A kernel-trace study of how container resource
limits slow down and break AI coding agents,'' arXiv, 2026.

\bibitem{ghrunners}
GitHub, ``GitHub-hosted runners reference,''
\url{https://docs.github.com/en/actions/reference/runners/github-hosted-runners}, accessed 2026-11-12.

\bibitem{e2bsandbox}
E2B, ``Sandbox customization: CPU and RAM,''
\url{https://e2b.dev/docs}, accessed 2026-11-12.

\bibitem{cloudrun}
Google Cloud, ``Configure memory and CPU limits for Cloud Run services,''
\url{https://cloud.google.com/run/docs/configuring/services/memory-limits}, accessed 2026-11-12.

\bibitem{fargate}
Amazon Web Services, ``Amazon ECS task CPU and memory for Fargate,''
\url{https://docs.aws.amazon.com/AmazonECS/latest/developerguide/fargate-tasks-services.html}, accessed 2026-11-12.

\bibitem{denning1978}
P.~J. Denning and J.~P. Buzen, ``The operational analysis of queueing
network models,'' \textit{ACM Computing Surveys}, vol.~10, no.~3,
pp.~225--261, 1978.

\end{thebibliography}


\end{document}
